% Pertemuan 3 Praktikum Pemrograman (IF25-11002). Dihasilkan oleh tools/build.py dari tools/konten/p03.py.
% Kompilasi: xelatex P03.tex (dua kali). Catatan: xelatex -jobname=P03-catatan "\def\catatan{1}% Pertemuan 3 Praktikum Pemrograman (IF25-11002). Dihasilkan oleh tools/build.py dari tools/konten/p03.py.
% Kompilasi: xelatex P03.tex (dua kali). Catatan: xelatex -jobname=P03-catatan "\def\catatan{1}% Pertemuan 3 Praktikum Pemrograman (IF25-11002). Dihasilkan oleh tools/build.py dari tools/konten/p03.py.
% Kompilasi: xelatex P03.tex (dua kali). Catatan: xelatex -jobname=P03-catatan "\def\catatan{1}% Pertemuan 3 Praktikum Pemrograman (IF25-11002). Dihasilkan oleh tools/build.py dari tools/konten/p03.py.
% Kompilasi: xelatex P03.tex (dua kali). Catatan: xelatex -jobname=P03-catatan "\def\catatan{1}\input{P03.tex}"
\documentclass[t]{beamer}
\usetheme{ITERA}
\input{catatan-preamble.tex}
\renewcommand\ITERAmeeting{Pertemuan 3}
\begin{document}
\SlideJudul{IF25-11002 PRAKTIKUM PEMROGRAMAN}{Pertemuan 3\\Operator dan Precedence}{Sisa bagi, urutan operasi, penugasan majemuk, serta operator relasional dan logika}{Tim Pengajar}{Tahun 2026. Sejajar dengan Pertemuan 3 Algoritma Pemrograman: Operator dan Precedence.}
\note{Buka dengan menanyakan satu hal dari pertemuan lalu yang masih membingungkan.}

\begin{frame}{Dua mata kuliah, satu minggu}
\framesubtitle{Sebelum mulai}
Topik minggu ini sama di dua mata kuliah, dengan peran yang berbeda.
\vspace{10pt}

\begin{columns}[T]
\begin{column}{0.47\textwidth}
\begin{Blok}{Algoritma: Operator dan Precedence}
Jenis operator, urutan pengerjaan (precedence), dan evaluasi ekspresi langkah demi langkah.
\end{Blok}
\end{column}
\begin{column}{0.47\textwidth}
\begin{BlokGelap}{Praktikum: Operator}
Menulis ekspresi C++ yang benar urutannya, memakai \texttt{\%}, \texttt{+=}, \texttt{++}, serta operator relasional dan logika.
\end{BlokGelap}
\end{column}
\end{columns}
\note{Tanyakan apa yang sudah dibahas di Algoritma minggu ini. Gunakan jawabannya sebagai jembatan ke praktikum.}
\end{frame}

\begin{frame}{Saat pulang nanti, kalian bisa}
\framesubtitle{Target hari ini}
\begin{columns}[T]
\begin{column}{0.47\textwidth}
\begin{Langkah}{1}{Menerapkan}
Menulis program yang memakai operator aritmetika termasuk \texttt{\%}, operator penugasan majemuk, increment, serta operator relasional dan logika untuk menyelesaikan perhitungan

{\footnotesize\color{ITERAInkMuted}Sub-CPMK 3.1, CPMK0607}
\end{Langkah}
\end{column}
\begin{column}{0.47\textwidth}
\begin{Langkah}{2}{Menjelaskan}
Menjelaskan urutan evaluasi sebuah ekspresi berdasarkan precedence dan asosiativitas, serta nilai yang dihasilkan pada setiap langkah

{\footnotesize\color{ITERAInkMuted}Sub-CPMK 3.2, CPMK0608}
\end{Langkah}
\end{column}
\end{columns}
\vspace{8pt}
\begin{Catatan}
Dua kemampuan ini dinilai sama berat di Exit Ticket, tugas, UTS, dan UAS.
\end{Catatan}
\note{Bacakan target dengan pelan. Kata kuncinya menerapkan dan menjelaskan.}
\end{frame}

\begin{frame}{Ingat minggu lalu}
\framesubtitle{Review, 5 menit}
\begin{columns}[T]
\begin{column}{0.31\textwidth}
\begin{BlokGelap}{Pertanyaan 1}
Tipe apa yang tepat untuk suhu 36.6?
\end{BlokGelap}
\end{column}
\begin{column}{0.31\textwidth}
\begin{BlokGelap}{Pertanyaan 2}
Berapa hasil \texttt{7 / 2} di C++, dan mengapa?
\end{BlokGelap}
\end{column}
\begin{column}{0.31\textwidth}
\begin{BlokGelap}{Pertanyaan 3}
Kapan perlu \texttt{cin.ignore()} sebelum \texttt{getline}?
\end{BlokGelap}
\end{column}
\end{columns}
\vspace{8pt}
Jawab di kertas dulu, lalu cocokkan dengan teman sebelah.\note{Pilih dua mahasiswa untuk menjawab. Koreksi singkat, jangan mengajar ulang.}
\end{frame}

\Bagian{1}{Masalah pembuka}{Durasi lagu di playlist}
\note{Masalah pembuka 10 menit.}

\begin{frame}[fragile]{Durasi lagu di playlist}
\framesubtitle{Masalah minggu ini}
\begin{columns}[T]
\begin{column}{0.55\textwidth}
{Aplikasi musik menyimpan durasi playlist dalam detik, misalnya 3725 detik. Pengguna ingin melihatnya sebagai jam, menit, dan detik.

\vspace{6pt}
Pembagian biasa memberi 1.03 jam, yang tidak membantu. Kita perlu cara mengambil bagian bulat dan sisanya.\par}
\vspace{6pt}
\begin{Catatan}
\textbf{Diskusi 2 menit.} Bagaimana kalian menghitung ini dengan tangan? Operasi apa yang dipakai untuk mendapat angka 5 di akhir?
\end{Catatan}
\end{column}
\begin{column}{0.41\textwidth}
\begin{Keluaran}[]
Durasi: 3725 detik
= 1 jam 2 menit 5 detik
\end{Keluaran}
\end{column}
\end{columns}
\note{Tampung beberapa jawaban tanpa menilai benar salah.}
\end{frame}

\begin{frame}{Dari langkah ke instruksi}
\framesubtitle{Sambungan dengan Algoritma}
\begin{columns}[T]
\begin{column}{0.31\textwidth}
\begin{Langkah}{1}{Pecah}
Ambil jam, lalu sisa detiknya, lalu menit, lalu sisa lagi.
\end{Langkah}
\end{column}
\begin{column}{0.31\textwidth}
\begin{Langkah}{2}{Rumuskan}
Jam adalah hasil bagi bulat; sisa adalah sisa bagi.
\end{Langkah}
\end{column}
\begin{column}{0.31\textwidth}
\begin{Langkah}{3}{Terjemahkan}
Pakai \texttt{/} untuk hasil bagi bulat dan \texttt{\%} untuk sisa.
\end{Langkah}
\end{column}
\end{columns}
\vspace{10pt}
Langkah 1 dan 2 dilatih di Algoritma. Langkah 3 kita kerjakan hari ini dengan C++.\note{Hubungkan eksplisit ke Algoritma minggu ini.}
\end{frame}

\Bagian{2}{Coba dulu, baru dijelaskan}{25 menit eksperimen di GDB Online}
\note{Struggle 25 menit. Berkeliling, jangan menjelaskan materi dulu.}

\begin{frame}{Misi kalian}
\framesubtitle{Struggle, 25 menit}
\begin{columns}[T]
\begin{column}{0.31\textwidth}
\begin{Langkah}{1}{Hitung}
Tampilkan \texttt{3725 / 3600} dan \texttt{3725 \% 3600}. Apa artinya?
\end{Langkah}
\end{column}
\begin{column}{0.31\textwidth}
\begin{Langkah}{2}{Urutkan}
Tebak lalu jalankan \texttt{cout << 2 + 3 * 4;} dan \texttt{cout << 10 - 4 - 3;}.
\end{Langkah}
\end{column}
\begin{column}{0.31\textwidth}
\begin{Langkah}{3}{Bandingkan}
Tampilkan \texttt{5 > 3} dan \texttt{5 == 3}. Apa yang muncul?
\end{Langkah}
\end{column}
\end{columns}
\vspace{8pt}
\begin{columns}[T]
\begin{column}{0.47\textwidth}
\begin{Blok}{Boleh}
Bertanya ke teman, mengubah apa saja, sengaja membuat error.
\end{Blok}
\end{column}
\begin{column}{0.47\textwidth}
\begin{BlokAlert}{Belum dulu}
Mencari jawaban jadi di internet atau AI.
\end{BlokAlert}
\end{column}
\end{columns}
\note{Catat kesalahan yang sering muncul untuk dibahas di debrief.}
\end{frame}

\begin{frame}{Apa yang kalian temukan?}
\framesubtitle{Debrief struggle}
\begin{columns}[T]
\begin{column}{0.31\textwidth}
\begin{BlokGelap}{Pertanyaan 1}
Apa arti hasil \texttt{3725 \% 3600}?
\end{BlokGelap}
\end{column}
\begin{column}{0.31\textwidth}
\begin{BlokGelap}{Pertanyaan 2}
Mengapa \texttt{2 + 3 * 4} bukan 20?
\end{BlokGelap}
\end{column}
\begin{column}{0.31\textwidth}
\begin{BlokGelap}{Pertanyaan 3}
Mengapa perbandingan tampil sebagai 1 atau 0?
\end{BlokGelap}
\end{column}
\end{columns}
\vspace{10pt}
Jawaban kalian akan kita uji di bagian berikutnya.\note{Tulis jawaban di papan; buktikan atau patahkan di bagian materi.}
\end{frame}

\Bagian{3}{Mengolah nilai dengan operator}{Aritmetika, urutan operasi, penugasan, perbandingan, dan logika}
\note{Materi 40 menit. Tunjukkan setiap contoh langsung di GDB Online.}

\begin{frame}[fragile]{Sisa bagi dengan \%}
\framesubtitle{Operator modulo}
\begin{columns}[T]
\begin{column}{0.60\textwidth}
\begin{Kode}[listing options={style=ITERACodeKecil}]{modulo.cpp}
int total = 3725;
int jam = total / 3600;
int menit = total % 3600 / 60;
int detik = total % 60;
cout << jam << " " << menit << " " << detik << endl;
cout << 17 % 5 << " " << 20 % 4 << endl;
\end{Kode}
\end{column}
\begin{column}{0.36\textwidth}
\only<1>{\begin{TebakDulu}
{\ITERAKickerFont\fontsize{10pt}{12pt}\selectfont\color{ITERAAlert}TEBAK DULU}\par\vspace{4pt}
Sebelum dijalankan, tulis keluarannya di kertas.
\end{TebakDulu}
}\only<2>{\KeluaranFile[listing options={style=ITERAPlainKecil}]{keluaran/P03/k001.txt}
\begin{Catatan}
\texttt{\%} hanya untuk bilangan bulat. \texttt{a \% b == 0} berarti a habis dibagi b.
\end{Catatan}
}
\end{column}
\end{columns}
\note<1>{Minta mahasiswa menebak keluaran sebelum membuka slide berikutnya. Tanyakan satu atau dua tebakan yang berbeda, lalu buktikan.}
\end{frame}

\begin{frame}{Urutan pengerjaan operator}
\framesubtitle{Precedence}
\small
\begin{tabular}{@{}>{\raggedright\arraybackslash}p{104pt}>{\raggedright\arraybackslash}p{274pt}>{\raggedright\arraybackslash}p{170pt}>{\raggedright\arraybackslash}p{302pt}@{}}
\kepala{Tingkat} & \kepala{Operator} & \kepala{Arah} & \kepala{Contoh}\\
1 (pertama) & \texttt{()} & dalam ke luar & \texttt{(2 + 3) * 4} = 20\\
\barisgenap 2 & \texttt{++ -- !} dan cast & kanan ke kiri & \texttt{!true} = false\\
3 & \texttt{* / \%} & kiri ke kanan & \texttt{20 / 4 * 2} = 10\\
\barisgenap 4 & \texttt{+ -} & kiri ke kanan & \texttt{10 - 4 - 3} = 3\\
5 & \texttt{< <= > >=} lalu \texttt{== !=} & kiri ke kanan & \texttt{3 < 5 == true}\\
\barisgenap 6 & \texttt{\&\&} lalu \texttt{||} & kiri ke kanan & \texttt{a || b \&\& c}\\
7 (terakhir) & \texttt{= += -= *= /= \%=} & kanan ke kiri & \texttt{a = b = 0}\\
\garisdasar
\end{tabular}
\vspace{10pt}

Ragu? Tambahkan kurung. Kurung tidak membuat program lambat, tetapi membuatnya mudah dibaca.
\end{frame}

\begin{frame}[fragile]{Menelusuri urutan operasi}
\framesubtitle{Evaluasi ekspresi}
\begin{columns}[T]
\begin{column}{0.55\textwidth}
\begin{Kode}[]{urutan.cpp}
cout << 2 + 3 * 4 << endl;
cout << (2 + 3) * 4 << endl;
cout << 20 / 4 * 2 << endl;
cout << 10 - 4 - 3 << endl;
cout << 7 + 10 % 4 * 2 << endl;
\end{Kode}
\end{column}
\begin{column}{0.41\textwidth}
\only<1>{\begin{TebakDulu}
{\ITERAKickerFont\fontsize{10pt}{12pt}\selectfont\color{ITERAAlert}TEBAK DULU}\par\vspace{4pt}
Sebelum dijalankan, tulis keluarannya di kertas.
\end{TebakDulu}
}\only<2>{\KeluaranFile[]{keluaran/P03/k002.txt}
\begin{Catatan}
Baris terakhir: \texttt{10 \% 4} = 2, lalu \texttt{2 * 2} = 4, lalu \texttt{7 + 4} = 11.
\end{Catatan}
}
\end{column}
\end{columns}
{\small Di Bagian B, tuliskan langkahnya satu per satu, bukan hanya hasil akhirnya.}
\note<1>{Kerjakan baris terakhir di papan langkah demi langkah dengan garis bawah pada operasi yang dikerjakan.}
\end{frame}

\begin{frame}[fragile]{Penugasan majemuk dan increment}
\framesubtitle{Menyingkat penugasan}
\begin{columns}[T]
\begin{column}{0.55\textwidth}
\begin{Kode}[]{majemuk.cpp}
int x = 5;
x += 3;
x *= 2;
x--;
int a = x++;
int b = ++x;
cout << a << " " << b << " " << x << endl;
\end{Kode}
\end{column}
\begin{column}{0.41\textwidth}
\only<1>{\begin{TebakDulu}
{\ITERAKickerFont\fontsize{10pt}{12pt}\selectfont\color{ITERAAlert}TEBAK DULU}\par\vspace{4pt}
Sebelum dijalankan, tulis keluarannya di kertas.
\end{TebakDulu}
}\only<2>{\KeluaranFile[]{keluaran/P03/k003.txt}
\begin{Catatan}
\texttt{x++} memakai nilai lama dulu, lalu menambah. \texttt{++x} menambah dulu, lalu memakai nilai baru.
\end{Catatan}
}
\end{column}
\end{columns}
\note<1>{Trace table: x 5, 8, 16, 15; a = 15 lalu x 16; x 17 lalu b = 17.}
\end{frame}

\begin{frame}[fragile]{Membandingkan dan menggabungkan kondisi}
\framesubtitle{Relasional dan logika}
\begin{columns}[T]
\begin{column}{0.55\textwidth}
\begin{Kode}[]{logika.cpp}
int nilai = 78, hadir = 12;
bool lulus = nilai >= 60;
bool aktif = hadir >= 11;
cout << boolalpha;
cout << lulus << " " << aktif << endl;
cout << (lulus && aktif) << endl;
cout << (nilai >= 85 || hadir == 14) << endl;
cout << !lulus << endl;
\end{Kode}
\end{column}
\begin{column}{0.41\textwidth}
\only<1>{\begin{TebakDulu}
{\ITERAKickerFont\fontsize{10pt}{12pt}\selectfont\color{ITERAAlert}TEBAK DULU}\par\vspace{4pt}
Sebelum dijalankan, tulis keluarannya di kertas.
\end{TebakDulu}
}\only<2>{\KeluaranFile[]{keluaran/P03/k004.txt}
\begin{Catatan}
\texttt{==} membandingkan, \texttt{=} menugaskan. Tanpa \texttt{boolalpha}, true dan false tampil sebagai 1 dan 0.
\end{Catatan}
}
\end{column}
\end{columns}
\note<1>{Minta mahasiswa menebak keluaran sebelum membuka slide berikutnya. Tanyakan satu atau dua tebakan yang berbeda, lalu buktikan.}
\end{frame}

\begin{frame}[fragile]{Konversi tipe dan fungsi matematika}
\framesubtitle{Cast dan cmath}
\begin{columns}[T]
\begin{column}{0.60\textwidth}
\begin{Kode}[listing options={style=ITERACodeKecil}]{konversi.cpp}
double d = 3.99;
int i = (int) d;
cout << i << endl;
cout << (double) 7 / 2 << endl;
cout << sqrt(16.0) << " " << pow(2, 10) << endl;
cout << round(2.5) << " " << floor(2.7) << endl;
\end{Kode}
\end{column}
\begin{column}{0.36\textwidth}
\only<1>{\begin{TebakDulu}
{\ITERAKickerFont\fontsize{10pt}{12pt}\selectfont\color{ITERAAlert}TEBAK DULU}\par\vspace{4pt}
Sebelum dijalankan, tulis keluarannya di kertas.
\end{TebakDulu}
}\only<2>{\KeluaranFile[listing options={style=ITERAPlainKecil}]{keluaran/P03/k005.txt}
\begin{Catatan}
Cast ke \texttt{int} memotong, tidak membulatkan. Fungsi matematika perlu \texttt{\#include <cmath>}.
\end{Catatan}
}
\end{column}
\end{columns}
\note<1>{Minta mahasiswa menebak keluaran sebelum membuka slide berikutnya. Tanyakan satu atau dua tebakan yang berbeda, lalu buktikan.}
\end{frame}

\begin{frame}[fragile]{Tebak keluarannya}
\framesubtitle{Latihan menjelaskan, 3 menit}
\begin{columns}[T]
\begin{column}{0.56\textwidth}
\begin{Kode}[]{tebak.cpp}
int a = 9, b = 4;
int c = a % b + a / b * 2;
a += c;
b = a-- - b;
cout << c << " " << a << " " << b << endl;
cout << (a > b && c != 5) << endl;
\end{Kode}
\end{column}
\begin{column}{0.40\textwidth}
\begin{Catatan}
Tulis keluarannya di kertas, karakter demi karakter. Jangan dijalankan dulu.
\end{Catatan}
\end{column}
\end{columns}
\note{Contoh soal Bagian B. Beri waktu 3 menit sebelum slide jawaban.}
\end{frame}

\begin{frame}[fragile]{Tebak keluarannya: jawaban}
\framesubtitle{Latihan menjelaskan}
\begin{columns}[T]
\begin{column}{0.56\textwidth}
\begin{Kode}[]{tebak.cpp}
int a = 9, b = 4;
int c = a % b + a / b * 2;
a += c;
b = a-- - b;
cout << c << " " << a << " " << b << endl;
cout << (a > b && c != 5) << endl;
\end{Kode}
\end{column}
\begin{column}{0.40\textwidth}
\begin{Keluaran}[]
5 13 10
0
\end{Keluaran}
\begin{Catatan}
\texttt{c} = 1 + 2 × 2 = 5. \texttt{a} menjadi 14. \texttt{b} dihitung dengan nilai lama \texttt{a} (14 − 4 = 10), lalu \texttt{a} turun ke 13. Karena \texttt{c} = 5, kondisi terakhir bernilai 0.
\end{Catatan}
\end{column}
\end{columns}
\note{Tanyakan jawaban yang paling sering salah, lalu telusuri bersama.}
\end{frame}

\begin{frame}{Error yang sering muncul minggu ini}
\framesubtitle{Pesan diambil langsung dari g++}
\footnotesize
\begin{tabular}{@{}>{\raggedright\arraybackslash}p{208pt}>{\raggedright\arraybackslash}p{256pt}>{\raggedright\arraybackslash}p{398pt}@{}}
\kepala{Kesalahan} & \kepala{Contoh} & \kepala{Pesan pertama dari g++}\\
\% pada bilangan pecahan & \texttt{5.5 \% 2} & \texttt{error: invalid operands of types 'double' and 'int' to binary 'operator\%'}\\
\barisgenap Ruas kiri bukan variabel & \texttt{5 = x;} & \texttt{error: lvalue required as left operand of assignment}\\
Pembagian dengan nol & \texttt{int y = 10 / 0;} & \texttt{warning: division by zero}\\
\barisgenap Perbandingan tanpa kurung di cout & \texttt{cout << 3 > 2;} & \texttt{error: no match for 'operator>'}\\
== yang tidak dipakai & \texttt{x == 5;} & \texttt{warning: statement has no effect}\\
\garisdasar
\end{tabular}
\vspace{10pt}

{\small Pesan di tabel ini dihasilkan dengan g++ dan opsi -Wall, sama seperti di cek.sh.}
\note{Minta mahasiswa memotret slide ini.}
\end{frame}

\Bagian{4}{Hands-on bertingkat}{45 menit, dari dasar sampai tantangan}
\note{Mahasiswa membuka folder 02 Hands-on/P3 - Operator - Hands-on.}

\begin{frame}{Peta hands-on}
\framesubtitle{Folder 02 Hands-on/P3 - Operator - Hands-on}
\small
\begin{tabular}{@{}>{\raggedright\arraybackslash}p{36pt}>{\raggedright\arraybackslash}p{267pt}>{\raggedright\arraybackslash}p{110pt}>{\raggedright\arraybackslash}p{83pt}>{\raggedright\arraybackslash}p{341pt}@{}}
\kepala{No} & \kepala{File} & \kepala{Tingkat} & \kepala{Waktu} & \kepala{Yang dilatih}\\
1 & \texttt{handson1-waktu.cpp} & Dasar & 10' & \texttt{/} dan \texttt{\%} untuk memecah satuan\\
\barisgenap 2 & \texttt{handson2-kembalian.cpp} & Menengah & 15' & \texttt{\%=} dan pembagian bulat berulang\\
3 & \texttt{handson3-rumus.cpp} & Tantangan & 20' & telusuri precedence, perbaiki tiga rumus\\
\barisgenap + & \texttt{bonus-genap.cpp} & Pengayaan & sisa & relasional, \texttt{\&\&}, \texttt{boolalpha}\\
\garisdasar
\end{tabular}
\vspace{10pt}

Kerjakan berurutan. Cocokkan keluaran dengan folder \texttt{expected/}; latihan yang membaca masukan memakai data di folder \texttt{input/}.\note{Saat berkeliling, minta mahasiswa yang mengerjakan latihan tantangan menjelaskan blok PENJELASAN mereka.}
\end{frame}

\begin{frame}{Cara mengecek dan aturannya}
\framesubtitle{Hands-on}
\begin{columns}[T]
\begin{column}{0.47\textwidth}
\begin{Blok}{Di GDB Online}
Salin isi file latihan, lengkapi TODO, klik Run. Bila program membaca masukan, ketik isi file di \texttt{input/}. Bandingkan dengan \texttt{expected/}.
\end{Blok}
\end{column}
\begin{column}{0.47\textwidth}
\begin{BlokGelap}{Di komputer sendiri}
Jalankan \texttt{bash cek.sh}. Skrip mengompilasi dengan \texttt{-Wall -Werror}, memberi masukan dari \texttt{input/}, lalu mencocokkan keluaran.
\end{BlokGelap}
\end{column}
\end{columns}
\vspace{6pt}
\begin{Catatan}
AI boleh untuk memahami pesan error, tidak untuk meminta solusi utuh. Exit Ticket dikerjakan tanpa bantuan apa pun.
\end{Catatan}
\note{Hands-on tidak dinilai langsung; yang dinilai Exit Ticket dan tugas.}
\end{frame}

\Bagian{5}{Exit Ticket dan tugas}{25 menit terakhir, lalu pekerjaan rumah}
\note{Mulai Exit Ticket tepat waktu.}

\begin{frame}{Exit Ticket 3}
\framesubtitle{Individual, 25 menit, tanpa AI dan internet}
\small
\begin{tabular}{@{}>{\raggedright\arraybackslash}p{60pt}>{\raggedright\arraybackslash}p{560pt}>{\raggedright\arraybackslash}p{80pt}>{\raggedright\arraybackslash}p{150pt}@{}}
\kepala{Soal} & \kepala{Isi} & \kepala{Poin} & \kepala{CPMK}\\
A1 & Tulis program yang memecah sebuah bilangan tiga digit menjadi ratusan, puluhan, dan satuan & 25 & CPMK0607\\
\barisgenap A2 & Tulis ekspresi logika untuk syarat yang diberikan dan tampilkan hasilnya & 25 & CPMK0607\\
B1 & Tunjukkan langkah evaluasi sebuah ekspresi campuran sesuai precedence & 20 & CPMK0608\\
\barisgenap B2 & Telusuri nilai variabel setelah penugasan majemuk dan increment & 20 & CPMK0608\\
B3 & Jelaskan beda \texttt{x++} dan \texttt{++x} dengan contoh & 10 & CPMK0608\\
\garisdasar
\end{tabular}
\vspace{10pt}

{\small Soal A dinilai dari keluaran yang tepat sama. Soal B dinilai dari ketepatan dan alasan. Pelanggaran aturan membuat nilai Exit Ticket ini 0.}
\note{Soal lengkap dibagikan terpisah dan tidak disimpan di repo publik.}
\end{frame}

\begin{frame}{Tugas 3: Kalkulator Lengkap dengan Validasi}
\framesubtitle{Dikumpulkan sebelum Pertemuan 4}
\begin{columns}[T]
\begin{column}{0.47\textwidth}
\begin{Blok}{Bagian A, program, 50 poin}
Buat program kalkulator yang membaca dua angka desimal, lalu menampilkan hasil berikut. Ketentuannya di slide berikut.
\end{Blok}
\end{column}
\begin{column}{0.47\textwidth}
\begin{BlokGelap}{Bagian B, penjelasan, 50 poin}
Komentar maksimal 150 kata: urutan evaluasi ekspresi syarat terakhir langkah demi langkah, alasan perlu cast sebelum \texttt{\%}, dan satu kesalahan precedence atau pembagian yang kalian temui.
\end{BlokGelap}
\end{column}
\end{columns}
\vspace{8pt}
{\small Data di tugas adalah data kalian sendiri. Rubrik lengkap ada di Modul Belajar 3.}
\note{Ingatkan bahwa penjelasan disalin dari AI mudah dikenali karena tidak menyebut pengalaman mereka sendiri.}
\end{frame}

\begin{frame}{Ketentuan Tugas 3}
\framesubtitle{Kalkulator Lengkap dengan Validasi}
\footnotesize
\begin{tabular}{@{}>{\raggedright\arraybackslash}p{87pt}>{\raggedright\arraybackslash}p{788pt}@{}}
\kepala{No} & \kepala{Ketentuan Bagian A}\\
1 & Masukan dua angka bertipe \texttt{double}\\
\barisgenap 2 & Hasil \texttt{+}, \texttt{-}, \texttt{*}, dan \texttt{/}\\
3 & Sisa bagi \texttt{\%} dari bagian bulat kedua angka (cast ke \texttt{int})\\
\barisgenap 4 & Sebelum membagi, periksa apakah pembagi nol; tampilkan pesan dengan operator ternary \texttt{?:}\\
5 & Nilai \texttt{bool}: apakah angka pertama lebih besar\\
\barisgenap 6 & Nilai \texttt{bool}: apakah hasil penjumlahan > 100 dan hasil kali > 50\\
Bonus & Tampilkan apakah bagian bulat masing-masing angka ganjil atau genap.\\
\garisdasar
\end{tabular}
\note{Bacakan ketentuan satu per satu dan tanyakan apakah ada yang belum jelas.}
\end{frame}

\begin{frame}{Rubrik Tugas 3}
\framesubtitle{Empat level, sama di semua instrumen}
\footnotesize
\begin{tabular}{@{}>{\raggedright\arraybackslash}p{166pt}>{\raggedright\arraybackslash}p{44pt}>{\raggedright\arraybackslash}p{166pt}>{\raggedright\arraybackslash}p{156pt}>{\raggedright\arraybackslash}p{146pt}>{\raggedright\arraybackslash}p{146pt}@{}}
\kepala{Kriteria} & \kepala{Poin} & \kepala{Sangat baik 85--100} & \kepala{Baik 70--84} & \kepala{Cukup 55--69} & \kepala{Kurang <55}\\
A1 Program berjalan & 20 & tanpa error dan peringatan & ada peringatan & perlu satu perbaikan & tidak terkompilasi\\
\barisgenap A2 Kelengkapan operasi & 15 & semua operasi dan validasi & satu operasi kurang & dua kurang & tiga atau lebih kurang\\
A3 Validasi dan logika & 15 & pembagi nol dan dua bool tepat & satu salah & dua salah & tidak ada validasi\\
\barisgenap B1 Langkah evaluasi & 30 & urutan tepat dan beralasan & satu langkah keliru & hanya hasil akhir & tidak ada atau disalin\\
B2 Cerita error dan pengujian & 20 & pesan, sebab, perbaikan, uji & pesan tidak dikutip & hanya perbaikan & tidak ada\\
\garisdasar
\end{tabular}
\vspace{8pt}

{\small Nilai kriteria = poin × skor level. Bagian A total 50, Bagian B total 50.}
\note{Rubrik ini sama strukturnya setiap minggu; yang berubah hanya isi kriteria A2, A3, dan B1.}
\end{frame}

\begin{frame}{Sebelum Pertemuan 4}
\framesubtitle{Persiapan}
\begin{columns}[T]
\begin{column}{0.31\textwidth}
\begin{Langkah}{1}{Selesaikan}
Hands-on yang belum lulus \texttt{cek.sh}, terutama latihan tantangan.
\end{Langkah}
\end{column}
\begin{column}{0.31\textwidth}
\begin{Langkah}{2}{Kumpulkan}
Tugas 3 sebelum Pertemuan 4 dimulai.
\end{Langkah}
\end{column}
\begin{column}{0.31\textwidth}
\begin{Langkah}{3}{Baca}
Modul Belajar 3 untuk mengulang, lalu bagian awal modul berikutnya.
\end{Langkah}
\end{column}
\end{columns}
\vspace{10pt}
Pertemuan 4 membahas percabangan dengan \texttt{if}, \texttt{else if}, dan \texttt{switch}, memakai ekspresi relasional dan logika dari minggu ini.\note{Tanyakan satu hal yang paling membingungkan hari ini untuk dibuka di review minggu depan.}
\end{frame}

\SlidePenutup{Terima kasih}{Sampai jumpa di Pertemuan 4: percabangan}
\note{Slide penutup.}
\end{document}
"
\documentclass[t]{beamer}
\usetheme{ITERA}
% Mode catatan pembicara: aktif bila \catatan didefinisikan sebelum \input.
\ifdefined\catatan
  \setbeameroption{show only notes}
  \setbeamertemplate{note page}{\vspace*{20pt}\hspace*{4pt}\begin{minipage}{900pt}
    {\ITERAKickerFont\fontsize{11pt}{14pt}\selectfont\color{ITERAGoldText}CATATAN PEMBICARA \ \ |\ \ SLIDE \insertframenumber}\par\vspace{4pt}
    {\ITERADisplay\bfseries\fontsize{22pt}{28pt}\selectfont\color{ITERAStructure}\insertframetitle}\par\vspace{12pt}
    \fontsize{15pt}{23pt}\selectfont\insertnote\end{minipage}}
\fi
\tikzset{fase/.style={draw=none,fill=#1,rounded corners=3pt,minimum width=158pt,minimum height=118pt,text width=146pt,align=center}}

\renewcommand\ITERAmeeting{Pertemuan 3}
\begin{document}
\SlideJudul{IF25-11002 PRAKTIKUM PEMROGRAMAN}{Pertemuan 3\\Operator dan Precedence}{Sisa bagi, urutan operasi, penugasan majemuk, serta operator relasional dan logika}{Tim Pengajar}{Tahun 2026. Sejajar dengan Pertemuan 3 Algoritma Pemrograman: Operator dan Precedence.}
\note{Buka dengan menanyakan satu hal dari pertemuan lalu yang masih membingungkan.}

\begin{frame}{Dua mata kuliah, satu minggu}
\framesubtitle{Sebelum mulai}
Topik minggu ini sama di dua mata kuliah, dengan peran yang berbeda.
\vspace{10pt}

\begin{columns}[T]
\begin{column}{0.47\textwidth}
\begin{Blok}{Algoritma: Operator dan Precedence}
Jenis operator, urutan pengerjaan (precedence), dan evaluasi ekspresi langkah demi langkah.
\end{Blok}
\end{column}
\begin{column}{0.47\textwidth}
\begin{BlokGelap}{Praktikum: Operator}
Menulis ekspresi C++ yang benar urutannya, memakai \texttt{\%}, \texttt{+=}, \texttt{++}, serta operator relasional dan logika.
\end{BlokGelap}
\end{column}
\end{columns}
\note{Tanyakan apa yang sudah dibahas di Algoritma minggu ini. Gunakan jawabannya sebagai jembatan ke praktikum.}
\end{frame}

\begin{frame}{Saat pulang nanti, kalian bisa}
\framesubtitle{Target hari ini}
\begin{columns}[T]
\begin{column}{0.47\textwidth}
\begin{Langkah}{1}{Menerapkan}
Menulis program yang memakai operator aritmetika termasuk \texttt{\%}, operator penugasan majemuk, increment, serta operator relasional dan logika untuk menyelesaikan perhitungan

{\footnotesize\color{ITERAInkMuted}Sub-CPMK 3.1, CPMK0607}
\end{Langkah}
\end{column}
\begin{column}{0.47\textwidth}
\begin{Langkah}{2}{Menjelaskan}
Menjelaskan urutan evaluasi sebuah ekspresi berdasarkan precedence dan asosiativitas, serta nilai yang dihasilkan pada setiap langkah

{\footnotesize\color{ITERAInkMuted}Sub-CPMK 3.2, CPMK0608}
\end{Langkah}
\end{column}
\end{columns}
\vspace{8pt}
\begin{Catatan}
Dua kemampuan ini dinilai sama berat di Exit Ticket, tugas, UTS, dan UAS.
\end{Catatan}
\note{Bacakan target dengan pelan. Kata kuncinya menerapkan dan menjelaskan.}
\end{frame}

\begin{frame}{Ingat minggu lalu}
\framesubtitle{Review, 5 menit}
\begin{columns}[T]
\begin{column}{0.31\textwidth}
\begin{BlokGelap}{Pertanyaan 1}
Tipe apa yang tepat untuk suhu 36.6?
\end{BlokGelap}
\end{column}
\begin{column}{0.31\textwidth}
\begin{BlokGelap}{Pertanyaan 2}
Berapa hasil \texttt{7 / 2} di C++, dan mengapa?
\end{BlokGelap}
\end{column}
\begin{column}{0.31\textwidth}
\begin{BlokGelap}{Pertanyaan 3}
Kapan perlu \texttt{cin.ignore()} sebelum \texttt{getline}?
\end{BlokGelap}
\end{column}
\end{columns}
\vspace{8pt}
Jawab di kertas dulu, lalu cocokkan dengan teman sebelah.\note{Pilih dua mahasiswa untuk menjawab. Koreksi singkat, jangan mengajar ulang.}
\end{frame}

\Bagian{1}{Masalah pembuka}{Durasi lagu di playlist}
\note{Masalah pembuka 10 menit.}

\begin{frame}[fragile]{Durasi lagu di playlist}
\framesubtitle{Masalah minggu ini}
\begin{columns}[T]
\begin{column}{0.55\textwidth}
{Aplikasi musik menyimpan durasi playlist dalam detik, misalnya 3725 detik. Pengguna ingin melihatnya sebagai jam, menit, dan detik.

\vspace{6pt}
Pembagian biasa memberi 1.03 jam, yang tidak membantu. Kita perlu cara mengambil bagian bulat dan sisanya.\par}
\vspace{6pt}
\begin{Catatan}
\textbf{Diskusi 2 menit.} Bagaimana kalian menghitung ini dengan tangan? Operasi apa yang dipakai untuk mendapat angka 5 di akhir?
\end{Catatan}
\end{column}
\begin{column}{0.41\textwidth}
\begin{Keluaran}[]
Durasi: 3725 detik
= 1 jam 2 menit 5 detik
\end{Keluaran}
\end{column}
\end{columns}
\note{Tampung beberapa jawaban tanpa menilai benar salah.}
\end{frame}

\begin{frame}{Dari langkah ke instruksi}
\framesubtitle{Sambungan dengan Algoritma}
\begin{columns}[T]
\begin{column}{0.31\textwidth}
\begin{Langkah}{1}{Pecah}
Ambil jam, lalu sisa detiknya, lalu menit, lalu sisa lagi.
\end{Langkah}
\end{column}
\begin{column}{0.31\textwidth}
\begin{Langkah}{2}{Rumuskan}
Jam adalah hasil bagi bulat; sisa adalah sisa bagi.
\end{Langkah}
\end{column}
\begin{column}{0.31\textwidth}
\begin{Langkah}{3}{Terjemahkan}
Pakai \texttt{/} untuk hasil bagi bulat dan \texttt{\%} untuk sisa.
\end{Langkah}
\end{column}
\end{columns}
\vspace{10pt}
Langkah 1 dan 2 dilatih di Algoritma. Langkah 3 kita kerjakan hari ini dengan C++.\note{Hubungkan eksplisit ke Algoritma minggu ini.}
\end{frame}

\Bagian{2}{Coba dulu, baru dijelaskan}{25 menit eksperimen di GDB Online}
\note{Struggle 25 menit. Berkeliling, jangan menjelaskan materi dulu.}

\begin{frame}{Misi kalian}
\framesubtitle{Struggle, 25 menit}
\begin{columns}[T]
\begin{column}{0.31\textwidth}
\begin{Langkah}{1}{Hitung}
Tampilkan \texttt{3725 / 3600} dan \texttt{3725 \% 3600}. Apa artinya?
\end{Langkah}
\end{column}
\begin{column}{0.31\textwidth}
\begin{Langkah}{2}{Urutkan}
Tebak lalu jalankan \texttt{cout << 2 + 3 * 4;} dan \texttt{cout << 10 - 4 - 3;}.
\end{Langkah}
\end{column}
\begin{column}{0.31\textwidth}
\begin{Langkah}{3}{Bandingkan}
Tampilkan \texttt{5 > 3} dan \texttt{5 == 3}. Apa yang muncul?
\end{Langkah}
\end{column}
\end{columns}
\vspace{8pt}
\begin{columns}[T]
\begin{column}{0.47\textwidth}
\begin{Blok}{Boleh}
Bertanya ke teman, mengubah apa saja, sengaja membuat error.
\end{Blok}
\end{column}
\begin{column}{0.47\textwidth}
\begin{BlokAlert}{Belum dulu}
Mencari jawaban jadi di internet atau AI.
\end{BlokAlert}
\end{column}
\end{columns}
\note{Catat kesalahan yang sering muncul untuk dibahas di debrief.}
\end{frame}

\begin{frame}{Apa yang kalian temukan?}
\framesubtitle{Debrief struggle}
\begin{columns}[T]
\begin{column}{0.31\textwidth}
\begin{BlokGelap}{Pertanyaan 1}
Apa arti hasil \texttt{3725 \% 3600}?
\end{BlokGelap}
\end{column}
\begin{column}{0.31\textwidth}
\begin{BlokGelap}{Pertanyaan 2}
Mengapa \texttt{2 + 3 * 4} bukan 20?
\end{BlokGelap}
\end{column}
\begin{column}{0.31\textwidth}
\begin{BlokGelap}{Pertanyaan 3}
Mengapa perbandingan tampil sebagai 1 atau 0?
\end{BlokGelap}
\end{column}
\end{columns}
\vspace{10pt}
Jawaban kalian akan kita uji di bagian berikutnya.\note{Tulis jawaban di papan; buktikan atau patahkan di bagian materi.}
\end{frame}

\Bagian{3}{Mengolah nilai dengan operator}{Aritmetika, urutan operasi, penugasan, perbandingan, dan logika}
\note{Materi 40 menit. Tunjukkan setiap contoh langsung di GDB Online.}

\begin{frame}[fragile]{Sisa bagi dengan \%}
\framesubtitle{Operator modulo}
\begin{columns}[T]
\begin{column}{0.60\textwidth}
\begin{Kode}[listing options={style=ITERACodeKecil}]{modulo.cpp}
int total = 3725;
int jam = total / 3600;
int menit = total % 3600 / 60;
int detik = total % 60;
cout << jam << " " << menit << " " << detik << endl;
cout << 17 % 5 << " " << 20 % 4 << endl;
\end{Kode}
\end{column}
\begin{column}{0.36\textwidth}
\only<1>{\begin{TebakDulu}
{\ITERAKickerFont\fontsize{10pt}{12pt}\selectfont\color{ITERAAlert}TEBAK DULU}\par\vspace{4pt}
Sebelum dijalankan, tulis keluarannya di kertas.
\end{TebakDulu}
}\only<2>{\KeluaranFile[listing options={style=ITERAPlainKecil}]{keluaran/P03/k001.txt}
\begin{Catatan}
\texttt{\%} hanya untuk bilangan bulat. \texttt{a \% b == 0} berarti a habis dibagi b.
\end{Catatan}
}
\end{column}
\end{columns}
\note<1>{Minta mahasiswa menebak keluaran sebelum membuka slide berikutnya. Tanyakan satu atau dua tebakan yang berbeda, lalu buktikan.}
\end{frame}

\begin{frame}{Urutan pengerjaan operator}
\framesubtitle{Precedence}
\small
\begin{tabular}{@{}>{\raggedright\arraybackslash}p{104pt}>{\raggedright\arraybackslash}p{274pt}>{\raggedright\arraybackslash}p{170pt}>{\raggedright\arraybackslash}p{302pt}@{}}
\kepala{Tingkat} & \kepala{Operator} & \kepala{Arah} & \kepala{Contoh}\\
1 (pertama) & \texttt{()} & dalam ke luar & \texttt{(2 + 3) * 4} = 20\\
\barisgenap 2 & \texttt{++ -- !} dan cast & kanan ke kiri & \texttt{!true} = false\\
3 & \texttt{* / \%} & kiri ke kanan & \texttt{20 / 4 * 2} = 10\\
\barisgenap 4 & \texttt{+ -} & kiri ke kanan & \texttt{10 - 4 - 3} = 3\\
5 & \texttt{< <= > >=} lalu \texttt{== !=} & kiri ke kanan & \texttt{3 < 5 == true}\\
\barisgenap 6 & \texttt{\&\&} lalu \texttt{||} & kiri ke kanan & \texttt{a || b \&\& c}\\
7 (terakhir) & \texttt{= += -= *= /= \%=} & kanan ke kiri & \texttt{a = b = 0}\\
\garisdasar
\end{tabular}
\vspace{10pt}

Ragu? Tambahkan kurung. Kurung tidak membuat program lambat, tetapi membuatnya mudah dibaca.
\end{frame}

\begin{frame}[fragile]{Menelusuri urutan operasi}
\framesubtitle{Evaluasi ekspresi}
\begin{columns}[T]
\begin{column}{0.55\textwidth}
\begin{Kode}[]{urutan.cpp}
cout << 2 + 3 * 4 << endl;
cout << (2 + 3) * 4 << endl;
cout << 20 / 4 * 2 << endl;
cout << 10 - 4 - 3 << endl;
cout << 7 + 10 % 4 * 2 << endl;
\end{Kode}
\end{column}
\begin{column}{0.41\textwidth}
\only<1>{\begin{TebakDulu}
{\ITERAKickerFont\fontsize{10pt}{12pt}\selectfont\color{ITERAAlert}TEBAK DULU}\par\vspace{4pt}
Sebelum dijalankan, tulis keluarannya di kertas.
\end{TebakDulu}
}\only<2>{\KeluaranFile[]{keluaran/P03/k002.txt}
\begin{Catatan}
Baris terakhir: \texttt{10 \% 4} = 2, lalu \texttt{2 * 2} = 4, lalu \texttt{7 + 4} = 11.
\end{Catatan}
}
\end{column}
\end{columns}
{\small Di Bagian B, tuliskan langkahnya satu per satu, bukan hanya hasil akhirnya.}
\note<1>{Kerjakan baris terakhir di papan langkah demi langkah dengan garis bawah pada operasi yang dikerjakan.}
\end{frame}

\begin{frame}[fragile]{Penugasan majemuk dan increment}
\framesubtitle{Menyingkat penugasan}
\begin{columns}[T]
\begin{column}{0.55\textwidth}
\begin{Kode}[]{majemuk.cpp}
int x = 5;
x += 3;
x *= 2;
x--;
int a = x++;
int b = ++x;
cout << a << " " << b << " " << x << endl;
\end{Kode}
\end{column}
\begin{column}{0.41\textwidth}
\only<1>{\begin{TebakDulu}
{\ITERAKickerFont\fontsize{10pt}{12pt}\selectfont\color{ITERAAlert}TEBAK DULU}\par\vspace{4pt}
Sebelum dijalankan, tulis keluarannya di kertas.
\end{TebakDulu}
}\only<2>{\KeluaranFile[]{keluaran/P03/k003.txt}
\begin{Catatan}
\texttt{x++} memakai nilai lama dulu, lalu menambah. \texttt{++x} menambah dulu, lalu memakai nilai baru.
\end{Catatan}
}
\end{column}
\end{columns}
\note<1>{Trace table: x 5, 8, 16, 15; a = 15 lalu x 16; x 17 lalu b = 17.}
\end{frame}

\begin{frame}[fragile]{Membandingkan dan menggabungkan kondisi}
\framesubtitle{Relasional dan logika}
\begin{columns}[T]
\begin{column}{0.55\textwidth}
\begin{Kode}[]{logika.cpp}
int nilai = 78, hadir = 12;
bool lulus = nilai >= 60;
bool aktif = hadir >= 11;
cout << boolalpha;
cout << lulus << " " << aktif << endl;
cout << (lulus && aktif) << endl;
cout << (nilai >= 85 || hadir == 14) << endl;
cout << !lulus << endl;
\end{Kode}
\end{column}
\begin{column}{0.41\textwidth}
\only<1>{\begin{TebakDulu}
{\ITERAKickerFont\fontsize{10pt}{12pt}\selectfont\color{ITERAAlert}TEBAK DULU}\par\vspace{4pt}
Sebelum dijalankan, tulis keluarannya di kertas.
\end{TebakDulu}
}\only<2>{\KeluaranFile[]{keluaran/P03/k004.txt}
\begin{Catatan}
\texttt{==} membandingkan, \texttt{=} menugaskan. Tanpa \texttt{boolalpha}, true dan false tampil sebagai 1 dan 0.
\end{Catatan}
}
\end{column}
\end{columns}
\note<1>{Minta mahasiswa menebak keluaran sebelum membuka slide berikutnya. Tanyakan satu atau dua tebakan yang berbeda, lalu buktikan.}
\end{frame}

\begin{frame}[fragile]{Konversi tipe dan fungsi matematika}
\framesubtitle{Cast dan cmath}
\begin{columns}[T]
\begin{column}{0.60\textwidth}
\begin{Kode}[listing options={style=ITERACodeKecil}]{konversi.cpp}
double d = 3.99;
int i = (int) d;
cout << i << endl;
cout << (double) 7 / 2 << endl;
cout << sqrt(16.0) << " " << pow(2, 10) << endl;
cout << round(2.5) << " " << floor(2.7) << endl;
\end{Kode}
\end{column}
\begin{column}{0.36\textwidth}
\only<1>{\begin{TebakDulu}
{\ITERAKickerFont\fontsize{10pt}{12pt}\selectfont\color{ITERAAlert}TEBAK DULU}\par\vspace{4pt}
Sebelum dijalankan, tulis keluarannya di kertas.
\end{TebakDulu}
}\only<2>{\KeluaranFile[listing options={style=ITERAPlainKecil}]{keluaran/P03/k005.txt}
\begin{Catatan}
Cast ke \texttt{int} memotong, tidak membulatkan. Fungsi matematika perlu \texttt{\#include <cmath>}.
\end{Catatan}
}
\end{column}
\end{columns}
\note<1>{Minta mahasiswa menebak keluaran sebelum membuka slide berikutnya. Tanyakan satu atau dua tebakan yang berbeda, lalu buktikan.}
\end{frame}

\begin{frame}[fragile]{Tebak keluarannya}
\framesubtitle{Latihan menjelaskan, 3 menit}
\begin{columns}[T]
\begin{column}{0.56\textwidth}
\begin{Kode}[]{tebak.cpp}
int a = 9, b = 4;
int c = a % b + a / b * 2;
a += c;
b = a-- - b;
cout << c << " " << a << " " << b << endl;
cout << (a > b && c != 5) << endl;
\end{Kode}
\end{column}
\begin{column}{0.40\textwidth}
\begin{Catatan}
Tulis keluarannya di kertas, karakter demi karakter. Jangan dijalankan dulu.
\end{Catatan}
\end{column}
\end{columns}
\note{Contoh soal Bagian B. Beri waktu 3 menit sebelum slide jawaban.}
\end{frame}

\begin{frame}[fragile]{Tebak keluarannya: jawaban}
\framesubtitle{Latihan menjelaskan}
\begin{columns}[T]
\begin{column}{0.56\textwidth}
\begin{Kode}[]{tebak.cpp}
int a = 9, b = 4;
int c = a % b + a / b * 2;
a += c;
b = a-- - b;
cout << c << " " << a << " " << b << endl;
cout << (a > b && c != 5) << endl;
\end{Kode}
\end{column}
\begin{column}{0.40\textwidth}
\begin{Keluaran}[]
5 13 10
0
\end{Keluaran}
\begin{Catatan}
\texttt{c} = 1 + 2 × 2 = 5. \texttt{a} menjadi 14. \texttt{b} dihitung dengan nilai lama \texttt{a} (14 − 4 = 10), lalu \texttt{a} turun ke 13. Karena \texttt{c} = 5, kondisi terakhir bernilai 0.
\end{Catatan}
\end{column}
\end{columns}
\note{Tanyakan jawaban yang paling sering salah, lalu telusuri bersama.}
\end{frame}

\begin{frame}{Error yang sering muncul minggu ini}
\framesubtitle{Pesan diambil langsung dari g++}
\footnotesize
\begin{tabular}{@{}>{\raggedright\arraybackslash}p{208pt}>{\raggedright\arraybackslash}p{256pt}>{\raggedright\arraybackslash}p{398pt}@{}}
\kepala{Kesalahan} & \kepala{Contoh} & \kepala{Pesan pertama dari g++}\\
\% pada bilangan pecahan & \texttt{5.5 \% 2} & \texttt{error: invalid operands of types 'double' and 'int' to binary 'operator\%'}\\
\barisgenap Ruas kiri bukan variabel & \texttt{5 = x;} & \texttt{error: lvalue required as left operand of assignment}\\
Pembagian dengan nol & \texttt{int y = 10 / 0;} & \texttt{warning: division by zero}\\
\barisgenap Perbandingan tanpa kurung di cout & \texttt{cout << 3 > 2;} & \texttt{error: no match for 'operator>'}\\
== yang tidak dipakai & \texttt{x == 5;} & \texttt{warning: statement has no effect}\\
\garisdasar
\end{tabular}
\vspace{10pt}

{\small Pesan di tabel ini dihasilkan dengan g++ dan opsi -Wall, sama seperti di cek.sh.}
\note{Minta mahasiswa memotret slide ini.}
\end{frame}

\Bagian{4}{Hands-on bertingkat}{45 menit, dari dasar sampai tantangan}
\note{Mahasiswa membuka folder 02 Hands-on/P3 - Operator - Hands-on.}

\begin{frame}{Peta hands-on}
\framesubtitle{Folder 02 Hands-on/P3 - Operator - Hands-on}
\small
\begin{tabular}{@{}>{\raggedright\arraybackslash}p{36pt}>{\raggedright\arraybackslash}p{267pt}>{\raggedright\arraybackslash}p{110pt}>{\raggedright\arraybackslash}p{83pt}>{\raggedright\arraybackslash}p{341pt}@{}}
\kepala{No} & \kepala{File} & \kepala{Tingkat} & \kepala{Waktu} & \kepala{Yang dilatih}\\
1 & \texttt{handson1-waktu.cpp} & Dasar & 10' & \texttt{/} dan \texttt{\%} untuk memecah satuan\\
\barisgenap 2 & \texttt{handson2-kembalian.cpp} & Menengah & 15' & \texttt{\%=} dan pembagian bulat berulang\\
3 & \texttt{handson3-rumus.cpp} & Tantangan & 20' & telusuri precedence, perbaiki tiga rumus\\
\barisgenap + & \texttt{bonus-genap.cpp} & Pengayaan & sisa & relasional, \texttt{\&\&}, \texttt{boolalpha}\\
\garisdasar
\end{tabular}
\vspace{10pt}

Kerjakan berurutan. Cocokkan keluaran dengan folder \texttt{expected/}; latihan yang membaca masukan memakai data di folder \texttt{input/}.\note{Saat berkeliling, minta mahasiswa yang mengerjakan latihan tantangan menjelaskan blok PENJELASAN mereka.}
\end{frame}

\begin{frame}{Cara mengecek dan aturannya}
\framesubtitle{Hands-on}
\begin{columns}[T]
\begin{column}{0.47\textwidth}
\begin{Blok}{Di GDB Online}
Salin isi file latihan, lengkapi TODO, klik Run. Bila program membaca masukan, ketik isi file di \texttt{input/}. Bandingkan dengan \texttt{expected/}.
\end{Blok}
\end{column}
\begin{column}{0.47\textwidth}
\begin{BlokGelap}{Di komputer sendiri}
Jalankan \texttt{bash cek.sh}. Skrip mengompilasi dengan \texttt{-Wall -Werror}, memberi masukan dari \texttt{input/}, lalu mencocokkan keluaran.
\end{BlokGelap}
\end{column}
\end{columns}
\vspace{6pt}
\begin{Catatan}
AI boleh untuk memahami pesan error, tidak untuk meminta solusi utuh. Exit Ticket dikerjakan tanpa bantuan apa pun.
\end{Catatan}
\note{Hands-on tidak dinilai langsung; yang dinilai Exit Ticket dan tugas.}
\end{frame}

\Bagian{5}{Exit Ticket dan tugas}{25 menit terakhir, lalu pekerjaan rumah}
\note{Mulai Exit Ticket tepat waktu.}

\begin{frame}{Exit Ticket 3}
\framesubtitle{Individual, 25 menit, tanpa AI dan internet}
\small
\begin{tabular}{@{}>{\raggedright\arraybackslash}p{60pt}>{\raggedright\arraybackslash}p{560pt}>{\raggedright\arraybackslash}p{80pt}>{\raggedright\arraybackslash}p{150pt}@{}}
\kepala{Soal} & \kepala{Isi} & \kepala{Poin} & \kepala{CPMK}\\
A1 & Tulis program yang memecah sebuah bilangan tiga digit menjadi ratusan, puluhan, dan satuan & 25 & CPMK0607\\
\barisgenap A2 & Tulis ekspresi logika untuk syarat yang diberikan dan tampilkan hasilnya & 25 & CPMK0607\\
B1 & Tunjukkan langkah evaluasi sebuah ekspresi campuran sesuai precedence & 20 & CPMK0608\\
\barisgenap B2 & Telusuri nilai variabel setelah penugasan majemuk dan increment & 20 & CPMK0608\\
B3 & Jelaskan beda \texttt{x++} dan \texttt{++x} dengan contoh & 10 & CPMK0608\\
\garisdasar
\end{tabular}
\vspace{10pt}

{\small Soal A dinilai dari keluaran yang tepat sama. Soal B dinilai dari ketepatan dan alasan. Pelanggaran aturan membuat nilai Exit Ticket ini 0.}
\note{Soal lengkap dibagikan terpisah dan tidak disimpan di repo publik.}
\end{frame}

\begin{frame}{Tugas 3: Kalkulator Lengkap dengan Validasi}
\framesubtitle{Dikumpulkan sebelum Pertemuan 4}
\begin{columns}[T]
\begin{column}{0.47\textwidth}
\begin{Blok}{Bagian A, program, 50 poin}
Buat program kalkulator yang membaca dua angka desimal, lalu menampilkan hasil berikut. Ketentuannya di slide berikut.
\end{Blok}
\end{column}
\begin{column}{0.47\textwidth}
\begin{BlokGelap}{Bagian B, penjelasan, 50 poin}
Komentar maksimal 150 kata: urutan evaluasi ekspresi syarat terakhir langkah demi langkah, alasan perlu cast sebelum \texttt{\%}, dan satu kesalahan precedence atau pembagian yang kalian temui.
\end{BlokGelap}
\end{column}
\end{columns}
\vspace{8pt}
{\small Data di tugas adalah data kalian sendiri. Rubrik lengkap ada di Modul Belajar 3.}
\note{Ingatkan bahwa penjelasan disalin dari AI mudah dikenali karena tidak menyebut pengalaman mereka sendiri.}
\end{frame}

\begin{frame}{Ketentuan Tugas 3}
\framesubtitle{Kalkulator Lengkap dengan Validasi}
\footnotesize
\begin{tabular}{@{}>{\raggedright\arraybackslash}p{87pt}>{\raggedright\arraybackslash}p{788pt}@{}}
\kepala{No} & \kepala{Ketentuan Bagian A}\\
1 & Masukan dua angka bertipe \texttt{double}\\
\barisgenap 2 & Hasil \texttt{+}, \texttt{-}, \texttt{*}, dan \texttt{/}\\
3 & Sisa bagi \texttt{\%} dari bagian bulat kedua angka (cast ke \texttt{int})\\
\barisgenap 4 & Sebelum membagi, periksa apakah pembagi nol; tampilkan pesan dengan operator ternary \texttt{?:}\\
5 & Nilai \texttt{bool}: apakah angka pertama lebih besar\\
\barisgenap 6 & Nilai \texttt{bool}: apakah hasil penjumlahan > 100 dan hasil kali > 50\\
Bonus & Tampilkan apakah bagian bulat masing-masing angka ganjil atau genap.\\
\garisdasar
\end{tabular}
\note{Bacakan ketentuan satu per satu dan tanyakan apakah ada yang belum jelas.}
\end{frame}

\begin{frame}{Rubrik Tugas 3}
\framesubtitle{Empat level, sama di semua instrumen}
\footnotesize
\begin{tabular}{@{}>{\raggedright\arraybackslash}p{166pt}>{\raggedright\arraybackslash}p{44pt}>{\raggedright\arraybackslash}p{166pt}>{\raggedright\arraybackslash}p{156pt}>{\raggedright\arraybackslash}p{146pt}>{\raggedright\arraybackslash}p{146pt}@{}}
\kepala{Kriteria} & \kepala{Poin} & \kepala{Sangat baik 85--100} & \kepala{Baik 70--84} & \kepala{Cukup 55--69} & \kepala{Kurang <55}\\
A1 Program berjalan & 20 & tanpa error dan peringatan & ada peringatan & perlu satu perbaikan & tidak terkompilasi\\
\barisgenap A2 Kelengkapan operasi & 15 & semua operasi dan validasi & satu operasi kurang & dua kurang & tiga atau lebih kurang\\
A3 Validasi dan logika & 15 & pembagi nol dan dua bool tepat & satu salah & dua salah & tidak ada validasi\\
\barisgenap B1 Langkah evaluasi & 30 & urutan tepat dan beralasan & satu langkah keliru & hanya hasil akhir & tidak ada atau disalin\\
B2 Cerita error dan pengujian & 20 & pesan, sebab, perbaikan, uji & pesan tidak dikutip & hanya perbaikan & tidak ada\\
\garisdasar
\end{tabular}
\vspace{8pt}

{\small Nilai kriteria = poin × skor level. Bagian A total 50, Bagian B total 50.}
\note{Rubrik ini sama strukturnya setiap minggu; yang berubah hanya isi kriteria A2, A3, dan B1.}
\end{frame}

\begin{frame}{Sebelum Pertemuan 4}
\framesubtitle{Persiapan}
\begin{columns}[T]
\begin{column}{0.31\textwidth}
\begin{Langkah}{1}{Selesaikan}
Hands-on yang belum lulus \texttt{cek.sh}, terutama latihan tantangan.
\end{Langkah}
\end{column}
\begin{column}{0.31\textwidth}
\begin{Langkah}{2}{Kumpulkan}
Tugas 3 sebelum Pertemuan 4 dimulai.
\end{Langkah}
\end{column}
\begin{column}{0.31\textwidth}
\begin{Langkah}{3}{Baca}
Modul Belajar 3 untuk mengulang, lalu bagian awal modul berikutnya.
\end{Langkah}
\end{column}
\end{columns}
\vspace{10pt}
Pertemuan 4 membahas percabangan dengan \texttt{if}, \texttt{else if}, dan \texttt{switch}, memakai ekspresi relasional dan logika dari minggu ini.\note{Tanyakan satu hal yang paling membingungkan hari ini untuk dibuka di review minggu depan.}
\end{frame}

\SlidePenutup{Terima kasih}{Sampai jumpa di Pertemuan 4: percabangan}
\note{Slide penutup.}
\end{document}
"
\documentclass[t]{beamer}
\usetheme{ITERA}
% Mode catatan pembicara: aktif bila \catatan didefinisikan sebelum \input.
\ifdefined\catatan
  \setbeameroption{show only notes}
  \setbeamertemplate{note page}{\vspace*{20pt}\hspace*{4pt}\begin{minipage}{900pt}
    {\ITERAKickerFont\fontsize{11pt}{14pt}\selectfont\color{ITERAGoldText}CATATAN PEMBICARA \ \ |\ \ SLIDE \insertframenumber}\par\vspace{4pt}
    {\ITERADisplay\bfseries\fontsize{22pt}{28pt}\selectfont\color{ITERAStructure}\insertframetitle}\par\vspace{12pt}
    \fontsize{15pt}{23pt}\selectfont\insertnote\end{minipage}}
\fi
\tikzset{fase/.style={draw=none,fill=#1,rounded corners=3pt,minimum width=158pt,minimum height=118pt,text width=146pt,align=center}}

\renewcommand\ITERAmeeting{Pertemuan 3}
\begin{document}
\SlideJudul{IF25-11002 PRAKTIKUM PEMROGRAMAN}{Pertemuan 3\\Operator dan Precedence}{Sisa bagi, urutan operasi, penugasan majemuk, serta operator relasional dan logika}{Tim Pengajar}{Tahun 2026. Sejajar dengan Pertemuan 3 Algoritma Pemrograman: Operator dan Precedence.}
\note{Buka dengan menanyakan satu hal dari pertemuan lalu yang masih membingungkan.}

\begin{frame}{Dua mata kuliah, satu minggu}
\framesubtitle{Sebelum mulai}
Topik minggu ini sama di dua mata kuliah, dengan peran yang berbeda.
\vspace{10pt}

\begin{columns}[T]
\begin{column}{0.47\textwidth}
\begin{Blok}{Algoritma: Operator dan Precedence}
Jenis operator, urutan pengerjaan (precedence), dan evaluasi ekspresi langkah demi langkah.
\end{Blok}
\end{column}
\begin{column}{0.47\textwidth}
\begin{BlokGelap}{Praktikum: Operator}
Menulis ekspresi C++ yang benar urutannya, memakai \texttt{\%}, \texttt{+=}, \texttt{++}, serta operator relasional dan logika.
\end{BlokGelap}
\end{column}
\end{columns}
\note{Tanyakan apa yang sudah dibahas di Algoritma minggu ini. Gunakan jawabannya sebagai jembatan ke praktikum.}
\end{frame}

\begin{frame}{Saat pulang nanti, kalian bisa}
\framesubtitle{Target hari ini}
\begin{columns}[T]
\begin{column}{0.47\textwidth}
\begin{Langkah}{1}{Menerapkan}
Menulis program yang memakai operator aritmetika termasuk \texttt{\%}, operator penugasan majemuk, increment, serta operator relasional dan logika untuk menyelesaikan perhitungan

{\footnotesize\color{ITERAInkMuted}Sub-CPMK 3.1, CPMK0607}
\end{Langkah}
\end{column}
\begin{column}{0.47\textwidth}
\begin{Langkah}{2}{Menjelaskan}
Menjelaskan urutan evaluasi sebuah ekspresi berdasarkan precedence dan asosiativitas, serta nilai yang dihasilkan pada setiap langkah

{\footnotesize\color{ITERAInkMuted}Sub-CPMK 3.2, CPMK0608}
\end{Langkah}
\end{column}
\end{columns}
\vspace{8pt}
\begin{Catatan}
Dua kemampuan ini dinilai sama berat di Exit Ticket, tugas, UTS, dan UAS.
\end{Catatan}
\note{Bacakan target dengan pelan. Kata kuncinya menerapkan dan menjelaskan.}
\end{frame}

\begin{frame}{Ingat minggu lalu}
\framesubtitle{Review, 5 menit}
\begin{columns}[T]
\begin{column}{0.31\textwidth}
\begin{BlokGelap}{Pertanyaan 1}
Tipe apa yang tepat untuk suhu 36.6?
\end{BlokGelap}
\end{column}
\begin{column}{0.31\textwidth}
\begin{BlokGelap}{Pertanyaan 2}
Berapa hasil \texttt{7 / 2} di C++, dan mengapa?
\end{BlokGelap}
\end{column}
\begin{column}{0.31\textwidth}
\begin{BlokGelap}{Pertanyaan 3}
Kapan perlu \texttt{cin.ignore()} sebelum \texttt{getline}?
\end{BlokGelap}
\end{column}
\end{columns}
\vspace{8pt}
Jawab di kertas dulu, lalu cocokkan dengan teman sebelah.\note{Pilih dua mahasiswa untuk menjawab. Koreksi singkat, jangan mengajar ulang.}
\end{frame}

\Bagian{1}{Masalah pembuka}{Durasi lagu di playlist}
\note{Masalah pembuka 10 menit.}

\begin{frame}[fragile]{Durasi lagu di playlist}
\framesubtitle{Masalah minggu ini}
\begin{columns}[T]
\begin{column}{0.55\textwidth}
{Aplikasi musik menyimpan durasi playlist dalam detik, misalnya 3725 detik. Pengguna ingin melihatnya sebagai jam, menit, dan detik.

\vspace{6pt}
Pembagian biasa memberi 1.03 jam, yang tidak membantu. Kita perlu cara mengambil bagian bulat dan sisanya.\par}
\vspace{6pt}
\begin{Catatan}
\textbf{Diskusi 2 menit.} Bagaimana kalian menghitung ini dengan tangan? Operasi apa yang dipakai untuk mendapat angka 5 di akhir?
\end{Catatan}
\end{column}
\begin{column}{0.41\textwidth}
\begin{Keluaran}[]
Durasi: 3725 detik
= 1 jam 2 menit 5 detik
\end{Keluaran}
\end{column}
\end{columns}
\note{Tampung beberapa jawaban tanpa menilai benar salah.}
\end{frame}

\begin{frame}{Dari langkah ke instruksi}
\framesubtitle{Sambungan dengan Algoritma}
\begin{columns}[T]
\begin{column}{0.31\textwidth}
\begin{Langkah}{1}{Pecah}
Ambil jam, lalu sisa detiknya, lalu menit, lalu sisa lagi.
\end{Langkah}
\end{column}
\begin{column}{0.31\textwidth}
\begin{Langkah}{2}{Rumuskan}
Jam adalah hasil bagi bulat; sisa adalah sisa bagi.
\end{Langkah}
\end{column}
\begin{column}{0.31\textwidth}
\begin{Langkah}{3}{Terjemahkan}
Pakai \texttt{/} untuk hasil bagi bulat dan \texttt{\%} untuk sisa.
\end{Langkah}
\end{column}
\end{columns}
\vspace{10pt}
Langkah 1 dan 2 dilatih di Algoritma. Langkah 3 kita kerjakan hari ini dengan C++.\note{Hubungkan eksplisit ke Algoritma minggu ini.}
\end{frame}

\Bagian{2}{Coba dulu, baru dijelaskan}{25 menit eksperimen di GDB Online}
\note{Struggle 25 menit. Berkeliling, jangan menjelaskan materi dulu.}

\begin{frame}{Misi kalian}
\framesubtitle{Struggle, 25 menit}
\begin{columns}[T]
\begin{column}{0.31\textwidth}
\begin{Langkah}{1}{Hitung}
Tampilkan \texttt{3725 / 3600} dan \texttt{3725 \% 3600}. Apa artinya?
\end{Langkah}
\end{column}
\begin{column}{0.31\textwidth}
\begin{Langkah}{2}{Urutkan}
Tebak lalu jalankan \texttt{cout << 2 + 3 * 4;} dan \texttt{cout << 10 - 4 - 3;}.
\end{Langkah}
\end{column}
\begin{column}{0.31\textwidth}
\begin{Langkah}{3}{Bandingkan}
Tampilkan \texttt{5 > 3} dan \texttt{5 == 3}. Apa yang muncul?
\end{Langkah}
\end{column}
\end{columns}
\vspace{8pt}
\begin{columns}[T]
\begin{column}{0.47\textwidth}
\begin{Blok}{Boleh}
Bertanya ke teman, mengubah apa saja, sengaja membuat error.
\end{Blok}
\end{column}
\begin{column}{0.47\textwidth}
\begin{BlokAlert}{Belum dulu}
Mencari jawaban jadi di internet atau AI.
\end{BlokAlert}
\end{column}
\end{columns}
\note{Catat kesalahan yang sering muncul untuk dibahas di debrief.}
\end{frame}

\begin{frame}{Apa yang kalian temukan?}
\framesubtitle{Debrief struggle}
\begin{columns}[T]
\begin{column}{0.31\textwidth}
\begin{BlokGelap}{Pertanyaan 1}
Apa arti hasil \texttt{3725 \% 3600}?
\end{BlokGelap}
\end{column}
\begin{column}{0.31\textwidth}
\begin{BlokGelap}{Pertanyaan 2}
Mengapa \texttt{2 + 3 * 4} bukan 20?
\end{BlokGelap}
\end{column}
\begin{column}{0.31\textwidth}
\begin{BlokGelap}{Pertanyaan 3}
Mengapa perbandingan tampil sebagai 1 atau 0?
\end{BlokGelap}
\end{column}
\end{columns}
\vspace{10pt}
Jawaban kalian akan kita uji di bagian berikutnya.\note{Tulis jawaban di papan; buktikan atau patahkan di bagian materi.}
\end{frame}

\Bagian{3}{Mengolah nilai dengan operator}{Aritmetika, urutan operasi, penugasan, perbandingan, dan logika}
\note{Materi 40 menit. Tunjukkan setiap contoh langsung di GDB Online.}

\begin{frame}[fragile]{Sisa bagi dengan \%}
\framesubtitle{Operator modulo}
\begin{columns}[T]
\begin{column}{0.60\textwidth}
\begin{Kode}[listing options={style=ITERACodeKecil}]{modulo.cpp}
int total = 3725;
int jam = total / 3600;
int menit = total % 3600 / 60;
int detik = total % 60;
cout << jam << " " << menit << " " << detik << endl;
cout << 17 % 5 << " " << 20 % 4 << endl;
\end{Kode}
\end{column}
\begin{column}{0.36\textwidth}
\only<1>{\begin{TebakDulu}
{\ITERAKickerFont\fontsize{10pt}{12pt}\selectfont\color{ITERAAlert}TEBAK DULU}\par\vspace{4pt}
Sebelum dijalankan, tulis keluarannya di kertas.
\end{TebakDulu}
}\only<2>{\KeluaranFile[listing options={style=ITERAPlainKecil}]{keluaran/P03/k001.txt}
\begin{Catatan}
\texttt{\%} hanya untuk bilangan bulat. \texttt{a \% b == 0} berarti a habis dibagi b.
\end{Catatan}
}
\end{column}
\end{columns}
\note<1>{Minta mahasiswa menebak keluaran sebelum membuka slide berikutnya. Tanyakan satu atau dua tebakan yang berbeda, lalu buktikan.}
\end{frame}

\begin{frame}{Urutan pengerjaan operator}
\framesubtitle{Precedence}
\small
\begin{tabular}{@{}>{\raggedright\arraybackslash}p{104pt}>{\raggedright\arraybackslash}p{274pt}>{\raggedright\arraybackslash}p{170pt}>{\raggedright\arraybackslash}p{302pt}@{}}
\kepala{Tingkat} & \kepala{Operator} & \kepala{Arah} & \kepala{Contoh}\\
1 (pertama) & \texttt{()} & dalam ke luar & \texttt{(2 + 3) * 4} = 20\\
\barisgenap 2 & \texttt{++ -- !} dan cast & kanan ke kiri & \texttt{!true} = false\\
3 & \texttt{* / \%} & kiri ke kanan & \texttt{20 / 4 * 2} = 10\\
\barisgenap 4 & \texttt{+ -} & kiri ke kanan & \texttt{10 - 4 - 3} = 3\\
5 & \texttt{< <= > >=} lalu \texttt{== !=} & kiri ke kanan & \texttt{3 < 5 == true}\\
\barisgenap 6 & \texttt{\&\&} lalu \texttt{||} & kiri ke kanan & \texttt{a || b \&\& c}\\
7 (terakhir) & \texttt{= += -= *= /= \%=} & kanan ke kiri & \texttt{a = b = 0}\\
\garisdasar
\end{tabular}
\vspace{10pt}

Ragu? Tambahkan kurung. Kurung tidak membuat program lambat, tetapi membuatnya mudah dibaca.
\end{frame}

\begin{frame}[fragile]{Menelusuri urutan operasi}
\framesubtitle{Evaluasi ekspresi}
\begin{columns}[T]
\begin{column}{0.55\textwidth}
\begin{Kode}[]{urutan.cpp}
cout << 2 + 3 * 4 << endl;
cout << (2 + 3) * 4 << endl;
cout << 20 / 4 * 2 << endl;
cout << 10 - 4 - 3 << endl;
cout << 7 + 10 % 4 * 2 << endl;
\end{Kode}
\end{column}
\begin{column}{0.41\textwidth}
\only<1>{\begin{TebakDulu}
{\ITERAKickerFont\fontsize{10pt}{12pt}\selectfont\color{ITERAAlert}TEBAK DULU}\par\vspace{4pt}
Sebelum dijalankan, tulis keluarannya di kertas.
\end{TebakDulu}
}\only<2>{\KeluaranFile[]{keluaran/P03/k002.txt}
\begin{Catatan}
Baris terakhir: \texttt{10 \% 4} = 2, lalu \texttt{2 * 2} = 4, lalu \texttt{7 + 4} = 11.
\end{Catatan}
}
\end{column}
\end{columns}
{\small Di Bagian B, tuliskan langkahnya satu per satu, bukan hanya hasil akhirnya.}
\note<1>{Kerjakan baris terakhir di papan langkah demi langkah dengan garis bawah pada operasi yang dikerjakan.}
\end{frame}

\begin{frame}[fragile]{Penugasan majemuk dan increment}
\framesubtitle{Menyingkat penugasan}
\begin{columns}[T]
\begin{column}{0.55\textwidth}
\begin{Kode}[]{majemuk.cpp}
int x = 5;
x += 3;
x *= 2;
x--;
int a = x++;
int b = ++x;
cout << a << " " << b << " " << x << endl;
\end{Kode}
\end{column}
\begin{column}{0.41\textwidth}
\only<1>{\begin{TebakDulu}
{\ITERAKickerFont\fontsize{10pt}{12pt}\selectfont\color{ITERAAlert}TEBAK DULU}\par\vspace{4pt}
Sebelum dijalankan, tulis keluarannya di kertas.
\end{TebakDulu}
}\only<2>{\KeluaranFile[]{keluaran/P03/k003.txt}
\begin{Catatan}
\texttt{x++} memakai nilai lama dulu, lalu menambah. \texttt{++x} menambah dulu, lalu memakai nilai baru.
\end{Catatan}
}
\end{column}
\end{columns}
\note<1>{Trace table: x 5, 8, 16, 15; a = 15 lalu x 16; x 17 lalu b = 17.}
\end{frame}

\begin{frame}[fragile]{Membandingkan dan menggabungkan kondisi}
\framesubtitle{Relasional dan logika}
\begin{columns}[T]
\begin{column}{0.55\textwidth}
\begin{Kode}[]{logika.cpp}
int nilai = 78, hadir = 12;
bool lulus = nilai >= 60;
bool aktif = hadir >= 11;
cout << boolalpha;
cout << lulus << " " << aktif << endl;
cout << (lulus && aktif) << endl;
cout << (nilai >= 85 || hadir == 14) << endl;
cout << !lulus << endl;
\end{Kode}
\end{column}
\begin{column}{0.41\textwidth}
\only<1>{\begin{TebakDulu}
{\ITERAKickerFont\fontsize{10pt}{12pt}\selectfont\color{ITERAAlert}TEBAK DULU}\par\vspace{4pt}
Sebelum dijalankan, tulis keluarannya di kertas.
\end{TebakDulu}
}\only<2>{\KeluaranFile[]{keluaran/P03/k004.txt}
\begin{Catatan}
\texttt{==} membandingkan, \texttt{=} menugaskan. Tanpa \texttt{boolalpha}, true dan false tampil sebagai 1 dan 0.
\end{Catatan}
}
\end{column}
\end{columns}
\note<1>{Minta mahasiswa menebak keluaran sebelum membuka slide berikutnya. Tanyakan satu atau dua tebakan yang berbeda, lalu buktikan.}
\end{frame}

\begin{frame}[fragile]{Konversi tipe dan fungsi matematika}
\framesubtitle{Cast dan cmath}
\begin{columns}[T]
\begin{column}{0.60\textwidth}
\begin{Kode}[listing options={style=ITERACodeKecil}]{konversi.cpp}
double d = 3.99;
int i = (int) d;
cout << i << endl;
cout << (double) 7 / 2 << endl;
cout << sqrt(16.0) << " " << pow(2, 10) << endl;
cout << round(2.5) << " " << floor(2.7) << endl;
\end{Kode}
\end{column}
\begin{column}{0.36\textwidth}
\only<1>{\begin{TebakDulu}
{\ITERAKickerFont\fontsize{10pt}{12pt}\selectfont\color{ITERAAlert}TEBAK DULU}\par\vspace{4pt}
Sebelum dijalankan, tulis keluarannya di kertas.
\end{TebakDulu}
}\only<2>{\KeluaranFile[listing options={style=ITERAPlainKecil}]{keluaran/P03/k005.txt}
\begin{Catatan}
Cast ke \texttt{int} memotong, tidak membulatkan. Fungsi matematika perlu \texttt{\#include <cmath>}.
\end{Catatan}
}
\end{column}
\end{columns}
\note<1>{Minta mahasiswa menebak keluaran sebelum membuka slide berikutnya. Tanyakan satu atau dua tebakan yang berbeda, lalu buktikan.}
\end{frame}

\begin{frame}[fragile]{Tebak keluarannya}
\framesubtitle{Latihan menjelaskan, 3 menit}
\begin{columns}[T]
\begin{column}{0.56\textwidth}
\begin{Kode}[]{tebak.cpp}
int a = 9, b = 4;
int c = a % b + a / b * 2;
a += c;
b = a-- - b;
cout << c << " " << a << " " << b << endl;
cout << (a > b && c != 5) << endl;
\end{Kode}
\end{column}
\begin{column}{0.40\textwidth}
\begin{Catatan}
Tulis keluarannya di kertas, karakter demi karakter. Jangan dijalankan dulu.
\end{Catatan}
\end{column}
\end{columns}
\note{Contoh soal Bagian B. Beri waktu 3 menit sebelum slide jawaban.}
\end{frame}

\begin{frame}[fragile]{Tebak keluarannya: jawaban}
\framesubtitle{Latihan menjelaskan}
\begin{columns}[T]
\begin{column}{0.56\textwidth}
\begin{Kode}[]{tebak.cpp}
int a = 9, b = 4;
int c = a % b + a / b * 2;
a += c;
b = a-- - b;
cout << c << " " << a << " " << b << endl;
cout << (a > b && c != 5) << endl;
\end{Kode}
\end{column}
\begin{column}{0.40\textwidth}
\begin{Keluaran}[]
5 13 10
0
\end{Keluaran}
\begin{Catatan}
\texttt{c} = 1 + 2 × 2 = 5. \texttt{a} menjadi 14. \texttt{b} dihitung dengan nilai lama \texttt{a} (14 − 4 = 10), lalu \texttt{a} turun ke 13. Karena \texttt{c} = 5, kondisi terakhir bernilai 0.
\end{Catatan}
\end{column}
\end{columns}
\note{Tanyakan jawaban yang paling sering salah, lalu telusuri bersama.}
\end{frame}

\begin{frame}{Error yang sering muncul minggu ini}
\framesubtitle{Pesan diambil langsung dari g++}
\footnotesize
\begin{tabular}{@{}>{\raggedright\arraybackslash}p{208pt}>{\raggedright\arraybackslash}p{256pt}>{\raggedright\arraybackslash}p{398pt}@{}}
\kepala{Kesalahan} & \kepala{Contoh} & \kepala{Pesan pertama dari g++}\\
\% pada bilangan pecahan & \texttt{5.5 \% 2} & \texttt{error: invalid operands of types 'double' and 'int' to binary 'operator\%'}\\
\barisgenap Ruas kiri bukan variabel & \texttt{5 = x;} & \texttt{error: lvalue required as left operand of assignment}\\
Pembagian dengan nol & \texttt{int y = 10 / 0;} & \texttt{warning: division by zero}\\
\barisgenap Perbandingan tanpa kurung di cout & \texttt{cout << 3 > 2;} & \texttt{error: no match for 'operator>'}\\
== yang tidak dipakai & \texttt{x == 5;} & \texttt{warning: statement has no effect}\\
\garisdasar
\end{tabular}
\vspace{10pt}

{\small Pesan di tabel ini dihasilkan dengan g++ dan opsi -Wall, sama seperti di cek.sh.}
\note{Minta mahasiswa memotret slide ini.}
\end{frame}

\Bagian{4}{Hands-on bertingkat}{45 menit, dari dasar sampai tantangan}
\note{Mahasiswa membuka folder 02 Hands-on/P3 - Operator - Hands-on.}

\begin{frame}{Peta hands-on}
\framesubtitle{Folder 02 Hands-on/P3 - Operator - Hands-on}
\small
\begin{tabular}{@{}>{\raggedright\arraybackslash}p{36pt}>{\raggedright\arraybackslash}p{267pt}>{\raggedright\arraybackslash}p{110pt}>{\raggedright\arraybackslash}p{83pt}>{\raggedright\arraybackslash}p{341pt}@{}}
\kepala{No} & \kepala{File} & \kepala{Tingkat} & \kepala{Waktu} & \kepala{Yang dilatih}\\
1 & \texttt{handson1-waktu.cpp} & Dasar & 10' & \texttt{/} dan \texttt{\%} untuk memecah satuan\\
\barisgenap 2 & \texttt{handson2-kembalian.cpp} & Menengah & 15' & \texttt{\%=} dan pembagian bulat berulang\\
3 & \texttt{handson3-rumus.cpp} & Tantangan & 20' & telusuri precedence, perbaiki tiga rumus\\
\barisgenap + & \texttt{bonus-genap.cpp} & Pengayaan & sisa & relasional, \texttt{\&\&}, \texttt{boolalpha}\\
\garisdasar
\end{tabular}
\vspace{10pt}

Kerjakan berurutan. Cocokkan keluaran dengan folder \texttt{expected/}; latihan yang membaca masukan memakai data di folder \texttt{input/}.\note{Saat berkeliling, minta mahasiswa yang mengerjakan latihan tantangan menjelaskan blok PENJELASAN mereka.}
\end{frame}

\begin{frame}{Cara mengecek dan aturannya}
\framesubtitle{Hands-on}
\begin{columns}[T]
\begin{column}{0.47\textwidth}
\begin{Blok}{Di GDB Online}
Salin isi file latihan, lengkapi TODO, klik Run. Bila program membaca masukan, ketik isi file di \texttt{input/}. Bandingkan dengan \texttt{expected/}.
\end{Blok}
\end{column}
\begin{column}{0.47\textwidth}
\begin{BlokGelap}{Di komputer sendiri}
Jalankan \texttt{bash cek.sh}. Skrip mengompilasi dengan \texttt{-Wall -Werror}, memberi masukan dari \texttt{input/}, lalu mencocokkan keluaran.
\end{BlokGelap}
\end{column}
\end{columns}
\vspace{6pt}
\begin{Catatan}
AI boleh untuk memahami pesan error, tidak untuk meminta solusi utuh. Exit Ticket dikerjakan tanpa bantuan apa pun.
\end{Catatan}
\note{Hands-on tidak dinilai langsung; yang dinilai Exit Ticket dan tugas.}
\end{frame}

\Bagian{5}{Exit Ticket dan tugas}{25 menit terakhir, lalu pekerjaan rumah}
\note{Mulai Exit Ticket tepat waktu.}

\begin{frame}{Exit Ticket 3}
\framesubtitle{Individual, 25 menit, tanpa AI dan internet}
\small
\begin{tabular}{@{}>{\raggedright\arraybackslash}p{60pt}>{\raggedright\arraybackslash}p{560pt}>{\raggedright\arraybackslash}p{80pt}>{\raggedright\arraybackslash}p{150pt}@{}}
\kepala{Soal} & \kepala{Isi} & \kepala{Poin} & \kepala{CPMK}\\
A1 & Tulis program yang memecah sebuah bilangan tiga digit menjadi ratusan, puluhan, dan satuan & 25 & CPMK0607\\
\barisgenap A2 & Tulis ekspresi logika untuk syarat yang diberikan dan tampilkan hasilnya & 25 & CPMK0607\\
B1 & Tunjukkan langkah evaluasi sebuah ekspresi campuran sesuai precedence & 20 & CPMK0608\\
\barisgenap B2 & Telusuri nilai variabel setelah penugasan majemuk dan increment & 20 & CPMK0608\\
B3 & Jelaskan beda \texttt{x++} dan \texttt{++x} dengan contoh & 10 & CPMK0608\\
\garisdasar
\end{tabular}
\vspace{10pt}

{\small Soal A dinilai dari keluaran yang tepat sama. Soal B dinilai dari ketepatan dan alasan. Pelanggaran aturan membuat nilai Exit Ticket ini 0.}
\note{Soal lengkap dibagikan terpisah dan tidak disimpan di repo publik.}
\end{frame}

\begin{frame}{Tugas 3: Kalkulator Lengkap dengan Validasi}
\framesubtitle{Dikumpulkan sebelum Pertemuan 4}
\begin{columns}[T]
\begin{column}{0.47\textwidth}
\begin{Blok}{Bagian A, program, 50 poin}
Buat program kalkulator yang membaca dua angka desimal, lalu menampilkan hasil berikut. Ketentuannya di slide berikut.
\end{Blok}
\end{column}
\begin{column}{0.47\textwidth}
\begin{BlokGelap}{Bagian B, penjelasan, 50 poin}
Komentar maksimal 150 kata: urutan evaluasi ekspresi syarat terakhir langkah demi langkah, alasan perlu cast sebelum \texttt{\%}, dan satu kesalahan precedence atau pembagian yang kalian temui.
\end{BlokGelap}
\end{column}
\end{columns}
\vspace{8pt}
{\small Data di tugas adalah data kalian sendiri. Rubrik lengkap ada di Modul Belajar 3.}
\note{Ingatkan bahwa penjelasan disalin dari AI mudah dikenali karena tidak menyebut pengalaman mereka sendiri.}
\end{frame}

\begin{frame}{Ketentuan Tugas 3}
\framesubtitle{Kalkulator Lengkap dengan Validasi}
\footnotesize
\begin{tabular}{@{}>{\raggedright\arraybackslash}p{87pt}>{\raggedright\arraybackslash}p{788pt}@{}}
\kepala{No} & \kepala{Ketentuan Bagian A}\\
1 & Masukan dua angka bertipe \texttt{double}\\
\barisgenap 2 & Hasil \texttt{+}, \texttt{-}, \texttt{*}, dan \texttt{/}\\
3 & Sisa bagi \texttt{\%} dari bagian bulat kedua angka (cast ke \texttt{int})\\
\barisgenap 4 & Sebelum membagi, periksa apakah pembagi nol; tampilkan pesan dengan operator ternary \texttt{?:}\\
5 & Nilai \texttt{bool}: apakah angka pertama lebih besar\\
\barisgenap 6 & Nilai \texttt{bool}: apakah hasil penjumlahan > 100 dan hasil kali > 50\\
Bonus & Tampilkan apakah bagian bulat masing-masing angka ganjil atau genap.\\
\garisdasar
\end{tabular}
\note{Bacakan ketentuan satu per satu dan tanyakan apakah ada yang belum jelas.}
\end{frame}

\begin{frame}{Rubrik Tugas 3}
\framesubtitle{Empat level, sama di semua instrumen}
\footnotesize
\begin{tabular}{@{}>{\raggedright\arraybackslash}p{166pt}>{\raggedright\arraybackslash}p{44pt}>{\raggedright\arraybackslash}p{166pt}>{\raggedright\arraybackslash}p{156pt}>{\raggedright\arraybackslash}p{146pt}>{\raggedright\arraybackslash}p{146pt}@{}}
\kepala{Kriteria} & \kepala{Poin} & \kepala{Sangat baik 85--100} & \kepala{Baik 70--84} & \kepala{Cukup 55--69} & \kepala{Kurang <55}\\
A1 Program berjalan & 20 & tanpa error dan peringatan & ada peringatan & perlu satu perbaikan & tidak terkompilasi\\
\barisgenap A2 Kelengkapan operasi & 15 & semua operasi dan validasi & satu operasi kurang & dua kurang & tiga atau lebih kurang\\
A3 Validasi dan logika & 15 & pembagi nol dan dua bool tepat & satu salah & dua salah & tidak ada validasi\\
\barisgenap B1 Langkah evaluasi & 30 & urutan tepat dan beralasan & satu langkah keliru & hanya hasil akhir & tidak ada atau disalin\\
B2 Cerita error dan pengujian & 20 & pesan, sebab, perbaikan, uji & pesan tidak dikutip & hanya perbaikan & tidak ada\\
\garisdasar
\end{tabular}
\vspace{8pt}

{\small Nilai kriteria = poin × skor level. Bagian A total 50, Bagian B total 50.}
\note{Rubrik ini sama strukturnya setiap minggu; yang berubah hanya isi kriteria A2, A3, dan B1.}
\end{frame}

\begin{frame}{Sebelum Pertemuan 4}
\framesubtitle{Persiapan}
\begin{columns}[T]
\begin{column}{0.31\textwidth}
\begin{Langkah}{1}{Selesaikan}
Hands-on yang belum lulus \texttt{cek.sh}, terutama latihan tantangan.
\end{Langkah}
\end{column}
\begin{column}{0.31\textwidth}
\begin{Langkah}{2}{Kumpulkan}
Tugas 3 sebelum Pertemuan 4 dimulai.
\end{Langkah}
\end{column}
\begin{column}{0.31\textwidth}
\begin{Langkah}{3}{Baca}
Modul Belajar 3 untuk mengulang, lalu bagian awal modul berikutnya.
\end{Langkah}
\end{column}
\end{columns}
\vspace{10pt}
Pertemuan 4 membahas percabangan dengan \texttt{if}, \texttt{else if}, dan \texttt{switch}, memakai ekspresi relasional dan logika dari minggu ini.\note{Tanyakan satu hal yang paling membingungkan hari ini untuk dibuka di review minggu depan.}
\end{frame}

\SlidePenutup{Terima kasih}{Sampai jumpa di Pertemuan 4: percabangan}
\note{Slide penutup.}
\end{document}
"
\documentclass[t]{beamer}
\usetheme{ITERA}
% Mode catatan pembicara: aktif bila \catatan didefinisikan sebelum \input.
\ifdefined\catatan
  \setbeameroption{show only notes}
  \setbeamertemplate{note page}{\vspace*{20pt}\hspace*{4pt}\begin{minipage}{900pt}
    {\ITERAKickerFont\fontsize{11pt}{14pt}\selectfont\color{ITERAGoldText}CATATAN PEMBICARA \ \ |\ \ SLIDE \insertframenumber}\par\vspace{4pt}
    {\ITERADisplay\bfseries\fontsize{22pt}{28pt}\selectfont\color{ITERAStructure}\insertframetitle}\par\vspace{12pt}
    \fontsize{15pt}{23pt}\selectfont\insertnote\end{minipage}}
\fi
\tikzset{fase/.style={draw=none,fill=#1,rounded corners=3pt,minimum width=158pt,minimum height=118pt,text width=146pt,align=center}}

\renewcommand\ITERAmeeting{Pertemuan 3}
\begin{document}
\SlideJudul{IF25-11002 PRAKTIKUM PEMROGRAMAN}{Pertemuan 3\\Operator dan Precedence}{Sisa bagi, urutan operasi, penugasan majemuk, serta operator relasional dan logika}{Tim Pengajar}{Tahun 2026. Sejajar dengan Pertemuan 3 Algoritma Pemrograman: Operator dan Precedence.}
\note{Buka dengan menanyakan satu hal dari pertemuan lalu yang masih membingungkan.}

\begin{frame}{Dua mata kuliah, satu minggu}
\framesubtitle{Sebelum mulai}
Topik minggu ini sama di dua mata kuliah, dengan peran yang berbeda.
\vspace{10pt}

\begin{columns}[T]
\begin{column}{0.47\textwidth}
\begin{Blok}{Algoritma: Operator dan Precedence}
Jenis operator, urutan pengerjaan (precedence), dan evaluasi ekspresi langkah demi langkah.
\end{Blok}
\end{column}
\begin{column}{0.47\textwidth}
\begin{BlokGelap}{Praktikum: Operator}
Menulis ekspresi C++ yang benar urutannya, memakai \texttt{\%}, \texttt{+=}, \texttt{++}, serta operator relasional dan logika.
\end{BlokGelap}
\end{column}
\end{columns}
\note{Tanyakan apa yang sudah dibahas di Algoritma minggu ini. Gunakan jawabannya sebagai jembatan ke praktikum.}
\end{frame}

\begin{frame}{Saat pulang nanti, kalian bisa}
\framesubtitle{Target hari ini}
\begin{columns}[T]
\begin{column}{0.47\textwidth}
\begin{Langkah}{1}{Menerapkan}
Menulis program yang memakai operator aritmetika termasuk \texttt{\%}, operator penugasan majemuk, increment, serta operator relasional dan logika untuk menyelesaikan perhitungan

{\footnotesize\color{ITERAInkMuted}Sub-CPMK 3.1, CPMK0607}
\end{Langkah}
\end{column}
\begin{column}{0.47\textwidth}
\begin{Langkah}{2}{Menjelaskan}
Menjelaskan urutan evaluasi sebuah ekspresi berdasarkan precedence dan asosiativitas, serta nilai yang dihasilkan pada setiap langkah

{\footnotesize\color{ITERAInkMuted}Sub-CPMK 3.2, CPMK0608}
\end{Langkah}
\end{column}
\end{columns}
\vspace{8pt}
\begin{Catatan}
Dua kemampuan ini dinilai sama berat di Exit Ticket, tugas, UTS, dan UAS.
\end{Catatan}
\note{Bacakan target dengan pelan. Kata kuncinya menerapkan dan menjelaskan.}
\end{frame}

\begin{frame}{Ingat minggu lalu}
\framesubtitle{Review, 5 menit}
\begin{columns}[T]
\begin{column}{0.31\textwidth}
\begin{BlokGelap}{Pertanyaan 1}
Tipe apa yang tepat untuk suhu 36.6?
\end{BlokGelap}
\end{column}
\begin{column}{0.31\textwidth}
\begin{BlokGelap}{Pertanyaan 2}
Berapa hasil \texttt{7 / 2} di C++, dan mengapa?
\end{BlokGelap}
\end{column}
\begin{column}{0.31\textwidth}
\begin{BlokGelap}{Pertanyaan 3}
Kapan perlu \texttt{cin.ignore()} sebelum \texttt{getline}?
\end{BlokGelap}
\end{column}
\end{columns}
\vspace{8pt}
Jawab di kertas dulu, lalu cocokkan dengan teman sebelah.\note{Pilih dua mahasiswa untuk menjawab. Koreksi singkat, jangan mengajar ulang.}
\end{frame}

\Bagian{1}{Masalah pembuka}{Durasi lagu di playlist}
\note{Masalah pembuka 10 menit.}

\begin{frame}[fragile]{Durasi lagu di playlist}
\framesubtitle{Masalah minggu ini}
\begin{columns}[T]
\begin{column}{0.55\textwidth}
{Aplikasi musik menyimpan durasi playlist dalam detik, misalnya 3725 detik. Pengguna ingin melihatnya sebagai jam, menit, dan detik.

\vspace{6pt}
Pembagian biasa memberi 1.03 jam, yang tidak membantu. Kita perlu cara mengambil bagian bulat dan sisanya.\par}
\vspace{6pt}
\begin{Catatan}
\textbf{Diskusi 2 menit.} Bagaimana kalian menghitung ini dengan tangan? Operasi apa yang dipakai untuk mendapat angka 5 di akhir?
\end{Catatan}
\end{column}
\begin{column}{0.41\textwidth}
\begin{Keluaran}[]
Durasi: 3725 detik
= 1 jam 2 menit 5 detik
\end{Keluaran}
\end{column}
\end{columns}
\note{Tampung beberapa jawaban tanpa menilai benar salah.}
\end{frame}

\begin{frame}{Dari langkah ke instruksi}
\framesubtitle{Sambungan dengan Algoritma}
\begin{columns}[T]
\begin{column}{0.31\textwidth}
\begin{Langkah}{1}{Pecah}
Ambil jam, lalu sisa detiknya, lalu menit, lalu sisa lagi.
\end{Langkah}
\end{column}
\begin{column}{0.31\textwidth}
\begin{Langkah}{2}{Rumuskan}
Jam adalah hasil bagi bulat; sisa adalah sisa bagi.
\end{Langkah}
\end{column}
\begin{column}{0.31\textwidth}
\begin{Langkah}{3}{Terjemahkan}
Pakai \texttt{/} untuk hasil bagi bulat dan \texttt{\%} untuk sisa.
\end{Langkah}
\end{column}
\end{columns}
\vspace{10pt}
Langkah 1 dan 2 dilatih di Algoritma. Langkah 3 kita kerjakan hari ini dengan C++.\note{Hubungkan eksplisit ke Algoritma minggu ini.}
\end{frame}

\Bagian{2}{Coba dulu, baru dijelaskan}{25 menit eksperimen di GDB Online}
\note{Struggle 25 menit. Berkeliling, jangan menjelaskan materi dulu.}

\begin{frame}{Misi kalian}
\framesubtitle{Struggle, 25 menit}
\begin{columns}[T]
\begin{column}{0.31\textwidth}
\begin{Langkah}{1}{Hitung}
Tampilkan \texttt{3725 / 3600} dan \texttt{3725 \% 3600}. Apa artinya?
\end{Langkah}
\end{column}
\begin{column}{0.31\textwidth}
\begin{Langkah}{2}{Urutkan}
Tebak lalu jalankan \texttt{cout << 2 + 3 * 4;} dan \texttt{cout << 10 - 4 - 3;}.
\end{Langkah}
\end{column}
\begin{column}{0.31\textwidth}
\begin{Langkah}{3}{Bandingkan}
Tampilkan \texttt{5 > 3} dan \texttt{5 == 3}. Apa yang muncul?
\end{Langkah}
\end{column}
\end{columns}
\vspace{8pt}
\begin{columns}[T]
\begin{column}{0.47\textwidth}
\begin{Blok}{Boleh}
Bertanya ke teman, mengubah apa saja, sengaja membuat error.
\end{Blok}
\end{column}
\begin{column}{0.47\textwidth}
\begin{BlokAlert}{Belum dulu}
Mencari jawaban jadi di internet atau AI.
\end{BlokAlert}
\end{column}
\end{columns}
\note{Catat kesalahan yang sering muncul untuk dibahas di debrief.}
\end{frame}

\begin{frame}{Apa yang kalian temukan?}
\framesubtitle{Debrief struggle}
\begin{columns}[T]
\begin{column}{0.31\textwidth}
\begin{BlokGelap}{Pertanyaan 1}
Apa arti hasil \texttt{3725 \% 3600}?
\end{BlokGelap}
\end{column}
\begin{column}{0.31\textwidth}
\begin{BlokGelap}{Pertanyaan 2}
Mengapa \texttt{2 + 3 * 4} bukan 20?
\end{BlokGelap}
\end{column}
\begin{column}{0.31\textwidth}
\begin{BlokGelap}{Pertanyaan 3}
Mengapa perbandingan tampil sebagai 1 atau 0?
\end{BlokGelap}
\end{column}
\end{columns}
\vspace{10pt}
Jawaban kalian akan kita uji di bagian berikutnya.\note{Tulis jawaban di papan; buktikan atau patahkan di bagian materi.}
\end{frame}

\Bagian{3}{Mengolah nilai dengan operator}{Aritmetika, urutan operasi, penugasan, perbandingan, dan logika}
\note{Materi 40 menit. Tunjukkan setiap contoh langsung di GDB Online.}

\begin{frame}[fragile]{Sisa bagi dengan \%}
\framesubtitle{Operator modulo}
\begin{columns}[T]
\begin{column}{0.60\textwidth}
\begin{Kode}[listing options={style=ITERACodeKecil}]{modulo.cpp}
int total = 3725;
int jam = total / 3600;
int menit = total % 3600 / 60;
int detik = total % 60;
cout << jam << " " << menit << " " << detik << endl;
cout << 17 % 5 << " " << 20 % 4 << endl;
\end{Kode}
\end{column}
\begin{column}{0.36\textwidth}
\only<1>{\begin{TebakDulu}
{\ITERAKickerFont\fontsize{10pt}{12pt}\selectfont\color{ITERAAlert}TEBAK DULU}\par\vspace{4pt}
Sebelum dijalankan, tulis keluarannya di kertas.
\end{TebakDulu}
}\only<2>{\KeluaranFile[listing options={style=ITERAPlainKecil}]{keluaran/P03/k001.txt}
\begin{Catatan}
\texttt{\%} hanya untuk bilangan bulat. \texttt{a \% b == 0} berarti a habis dibagi b.
\end{Catatan}
}
\end{column}
\end{columns}
\note<1>{Minta mahasiswa menebak keluaran sebelum membuka slide berikutnya. Tanyakan satu atau dua tebakan yang berbeda, lalu buktikan.}
\end{frame}

\begin{frame}{Urutan pengerjaan operator}
\framesubtitle{Precedence}
\small
\begin{tabular}{@{}>{\raggedright\arraybackslash}p{104pt}>{\raggedright\arraybackslash}p{274pt}>{\raggedright\arraybackslash}p{170pt}>{\raggedright\arraybackslash}p{302pt}@{}}
\kepala{Tingkat} & \kepala{Operator} & \kepala{Arah} & \kepala{Contoh}\\
1 (pertama) & \texttt{()} & dalam ke luar & \texttt{(2 + 3) * 4} = 20\\
\barisgenap 2 & \texttt{++ -- !} dan cast & kanan ke kiri & \texttt{!true} = false\\
3 & \texttt{* / \%} & kiri ke kanan & \texttt{20 / 4 * 2} = 10\\
\barisgenap 4 & \texttt{+ -} & kiri ke kanan & \texttt{10 - 4 - 3} = 3\\
5 & \texttt{< <= > >=} lalu \texttt{== !=} & kiri ke kanan & \texttt{3 < 5 == true}\\
\barisgenap 6 & \texttt{\&\&} lalu \texttt{||} & kiri ke kanan & \texttt{a || b \&\& c}\\
7 (terakhir) & \texttt{= += -= *= /= \%=} & kanan ke kiri & \texttt{a = b = 0}\\
\garisdasar
\end{tabular}
\vspace{10pt}

Ragu? Tambahkan kurung. Kurung tidak membuat program lambat, tetapi membuatnya mudah dibaca.
\end{frame}

\begin{frame}[fragile]{Menelusuri urutan operasi}
\framesubtitle{Evaluasi ekspresi}
\begin{columns}[T]
\begin{column}{0.55\textwidth}
\begin{Kode}[]{urutan.cpp}
cout << 2 + 3 * 4 << endl;
cout << (2 + 3) * 4 << endl;
cout << 20 / 4 * 2 << endl;
cout << 10 - 4 - 3 << endl;
cout << 7 + 10 % 4 * 2 << endl;
\end{Kode}
\end{column}
\begin{column}{0.41\textwidth}
\only<1>{\begin{TebakDulu}
{\ITERAKickerFont\fontsize{10pt}{12pt}\selectfont\color{ITERAAlert}TEBAK DULU}\par\vspace{4pt}
Sebelum dijalankan, tulis keluarannya di kertas.
\end{TebakDulu}
}\only<2>{\KeluaranFile[]{keluaran/P03/k002.txt}
\begin{Catatan}
Baris terakhir: \texttt{10 \% 4} = 2, lalu \texttt{2 * 2} = 4, lalu \texttt{7 + 4} = 11.
\end{Catatan}
}
\end{column}
\end{columns}
{\small Di Bagian B, tuliskan langkahnya satu per satu, bukan hanya hasil akhirnya.}
\note<1>{Kerjakan baris terakhir di papan langkah demi langkah dengan garis bawah pada operasi yang dikerjakan.}
\end{frame}

\begin{frame}[fragile]{Penugasan majemuk dan increment}
\framesubtitle{Menyingkat penugasan}
\begin{columns}[T]
\begin{column}{0.55\textwidth}
\begin{Kode}[]{majemuk.cpp}
int x = 5;
x += 3;
x *= 2;
x--;
int a = x++;
int b = ++x;
cout << a << " " << b << " " << x << endl;
\end{Kode}
\end{column}
\begin{column}{0.41\textwidth}
\only<1>{\begin{TebakDulu}
{\ITERAKickerFont\fontsize{10pt}{12pt}\selectfont\color{ITERAAlert}TEBAK DULU}\par\vspace{4pt}
Sebelum dijalankan, tulis keluarannya di kertas.
\end{TebakDulu}
}\only<2>{\KeluaranFile[]{keluaran/P03/k003.txt}
\begin{Catatan}
\texttt{x++} memakai nilai lama dulu, lalu menambah. \texttt{++x} menambah dulu, lalu memakai nilai baru.
\end{Catatan}
}
\end{column}
\end{columns}
\note<1>{Trace table: x 5, 8, 16, 15; a = 15 lalu x 16; x 17 lalu b = 17.}
\end{frame}

\begin{frame}[fragile]{Membandingkan dan menggabungkan kondisi}
\framesubtitle{Relasional dan logika}
\begin{columns}[T]
\begin{column}{0.55\textwidth}
\begin{Kode}[]{logika.cpp}
int nilai = 78, hadir = 12;
bool lulus = nilai >= 60;
bool aktif = hadir >= 11;
cout << boolalpha;
cout << lulus << " " << aktif << endl;
cout << (lulus && aktif) << endl;
cout << (nilai >= 85 || hadir == 14) << endl;
cout << !lulus << endl;
\end{Kode}
\end{column}
\begin{column}{0.41\textwidth}
\only<1>{\begin{TebakDulu}
{\ITERAKickerFont\fontsize{10pt}{12pt}\selectfont\color{ITERAAlert}TEBAK DULU}\par\vspace{4pt}
Sebelum dijalankan, tulis keluarannya di kertas.
\end{TebakDulu}
}\only<2>{\KeluaranFile[]{keluaran/P03/k004.txt}
\begin{Catatan}
\texttt{==} membandingkan, \texttt{=} menugaskan. Tanpa \texttt{boolalpha}, true dan false tampil sebagai 1 dan 0.
\end{Catatan}
}
\end{column}
\end{columns}
\note<1>{Minta mahasiswa menebak keluaran sebelum membuka slide berikutnya. Tanyakan satu atau dua tebakan yang berbeda, lalu buktikan.}
\end{frame}

\begin{frame}[fragile]{Konversi tipe dan fungsi matematika}
\framesubtitle{Cast dan cmath}
\begin{columns}[T]
\begin{column}{0.60\textwidth}
\begin{Kode}[listing options={style=ITERACodeKecil}]{konversi.cpp}
double d = 3.99;
int i = (int) d;
cout << i << endl;
cout << (double) 7 / 2 << endl;
cout << sqrt(16.0) << " " << pow(2, 10) << endl;
cout << round(2.5) << " " << floor(2.7) << endl;
\end{Kode}
\end{column}
\begin{column}{0.36\textwidth}
\only<1>{\begin{TebakDulu}
{\ITERAKickerFont\fontsize{10pt}{12pt}\selectfont\color{ITERAAlert}TEBAK DULU}\par\vspace{4pt}
Sebelum dijalankan, tulis keluarannya di kertas.
\end{TebakDulu}
}\only<2>{\KeluaranFile[listing options={style=ITERAPlainKecil}]{keluaran/P03/k005.txt}
\begin{Catatan}
Cast ke \texttt{int} memotong, tidak membulatkan. Fungsi matematika perlu \texttt{\#include <cmath>}.
\end{Catatan}
}
\end{column}
\end{columns}
\note<1>{Minta mahasiswa menebak keluaran sebelum membuka slide berikutnya. Tanyakan satu atau dua tebakan yang berbeda, lalu buktikan.}
\end{frame}

\begin{frame}[fragile]{Tebak keluarannya}
\framesubtitle{Latihan menjelaskan, 3 menit}
\begin{columns}[T]
\begin{column}{0.56\textwidth}
\begin{Kode}[]{tebak.cpp}
int a = 9, b = 4;
int c = a % b + a / b * 2;
a += c;
b = a-- - b;
cout << c << " " << a << " " << b << endl;
cout << (a > b && c != 5) << endl;
\end{Kode}
\end{column}
\begin{column}{0.40\textwidth}
\begin{Catatan}
Tulis keluarannya di kertas, karakter demi karakter. Jangan dijalankan dulu.
\end{Catatan}
\end{column}
\end{columns}
\note{Contoh soal Bagian B. Beri waktu 3 menit sebelum slide jawaban.}
\end{frame}

\begin{frame}[fragile]{Tebak keluarannya: jawaban}
\framesubtitle{Latihan menjelaskan}
\begin{columns}[T]
\begin{column}{0.56\textwidth}
\begin{Kode}[]{tebak.cpp}
int a = 9, b = 4;
int c = a % b + a / b * 2;
a += c;
b = a-- - b;
cout << c << " " << a << " " << b << endl;
cout << (a > b && c != 5) << endl;
\end{Kode}
\end{column}
\begin{column}{0.40\textwidth}
\begin{Keluaran}[]
5 13 10
0
\end{Keluaran}
\begin{Catatan}
\texttt{c} = 1 + 2 × 2 = 5. \texttt{a} menjadi 14. \texttt{b} dihitung dengan nilai lama \texttt{a} (14 − 4 = 10), lalu \texttt{a} turun ke 13. Karena \texttt{c} = 5, kondisi terakhir bernilai 0.
\end{Catatan}
\end{column}
\end{columns}
\note{Tanyakan jawaban yang paling sering salah, lalu telusuri bersama.}
\end{frame}

\begin{frame}{Error yang sering muncul minggu ini}
\framesubtitle{Pesan diambil langsung dari g++}
\footnotesize
\begin{tabular}{@{}>{\raggedright\arraybackslash}p{208pt}>{\raggedright\arraybackslash}p{256pt}>{\raggedright\arraybackslash}p{398pt}@{}}
\kepala{Kesalahan} & \kepala{Contoh} & \kepala{Pesan pertama dari g++}\\
\% pada bilangan pecahan & \texttt{5.5 \% 2} & \texttt{error: invalid operands of types 'double' and 'int' to binary 'operator\%'}\\
\barisgenap Ruas kiri bukan variabel & \texttt{5 = x;} & \texttt{error: lvalue required as left operand of assignment}\\
Pembagian dengan nol & \texttt{int y = 10 / 0;} & \texttt{warning: division by zero}\\
\barisgenap Perbandingan tanpa kurung di cout & \texttt{cout << 3 > 2;} & \texttt{error: no match for 'operator>'}\\
== yang tidak dipakai & \texttt{x == 5;} & \texttt{warning: statement has no effect}\\
\garisdasar
\end{tabular}
\vspace{10pt}

{\small Pesan di tabel ini dihasilkan dengan g++ dan opsi -Wall, sama seperti di cek.sh.}
\note{Minta mahasiswa memotret slide ini.}
\end{frame}

\Bagian{4}{Hands-on bertingkat}{45 menit, dari dasar sampai tantangan}
\note{Mahasiswa membuka folder 02 Hands-on/P3 - Operator - Hands-on.}

\begin{frame}{Peta hands-on}
\framesubtitle{Folder 02 Hands-on/P3 - Operator - Hands-on}
\small
\begin{tabular}{@{}>{\raggedright\arraybackslash}p{36pt}>{\raggedright\arraybackslash}p{267pt}>{\raggedright\arraybackslash}p{110pt}>{\raggedright\arraybackslash}p{83pt}>{\raggedright\arraybackslash}p{341pt}@{}}
\kepala{No} & \kepala{File} & \kepala{Tingkat} & \kepala{Waktu} & \kepala{Yang dilatih}\\
1 & \texttt{handson1-waktu.cpp} & Dasar & 10' & \texttt{/} dan \texttt{\%} untuk memecah satuan\\
\barisgenap 2 & \texttt{handson2-kembalian.cpp} & Menengah & 15' & \texttt{\%=} dan pembagian bulat berulang\\
3 & \texttt{handson3-rumus.cpp} & Tantangan & 20' & telusuri precedence, perbaiki tiga rumus\\
\barisgenap + & \texttt{bonus-genap.cpp} & Pengayaan & sisa & relasional, \texttt{\&\&}, \texttt{boolalpha}\\
\garisdasar
\end{tabular}
\vspace{10pt}

Kerjakan berurutan. Cocokkan keluaran dengan folder \texttt{expected/}; latihan yang membaca masukan memakai data di folder \texttt{input/}.\note{Saat berkeliling, minta mahasiswa yang mengerjakan latihan tantangan menjelaskan blok PENJELASAN mereka.}
\end{frame}

\begin{frame}{Cara mengecek dan aturannya}
\framesubtitle{Hands-on}
\begin{columns}[T]
\begin{column}{0.47\textwidth}
\begin{Blok}{Di GDB Online}
Salin isi file latihan, lengkapi TODO, klik Run. Bila program membaca masukan, ketik isi file di \texttt{input/}. Bandingkan dengan \texttt{expected/}.
\end{Blok}
\end{column}
\begin{column}{0.47\textwidth}
\begin{BlokGelap}{Di komputer sendiri}
Jalankan \texttt{bash cek.sh}. Skrip mengompilasi dengan \texttt{-Wall -Werror}, memberi masukan dari \texttt{input/}, lalu mencocokkan keluaran.
\end{BlokGelap}
\end{column}
\end{columns}
\vspace{6pt}
\begin{Catatan}
AI boleh untuk memahami pesan error, tidak untuk meminta solusi utuh. Exit Ticket dikerjakan tanpa bantuan apa pun.
\end{Catatan}
\note{Hands-on tidak dinilai langsung; yang dinilai Exit Ticket dan tugas.}
\end{frame}

\Bagian{5}{Exit Ticket dan tugas}{25 menit terakhir, lalu pekerjaan rumah}
\note{Mulai Exit Ticket tepat waktu.}

\begin{frame}{Exit Ticket 3}
\framesubtitle{Individual, 25 menit, tanpa AI dan internet}
\small
\begin{tabular}{@{}>{\raggedright\arraybackslash}p{60pt}>{\raggedright\arraybackslash}p{560pt}>{\raggedright\arraybackslash}p{80pt}>{\raggedright\arraybackslash}p{150pt}@{}}
\kepala{Soal} & \kepala{Isi} & \kepala{Poin} & \kepala{CPMK}\\
A1 & Tulis program yang memecah sebuah bilangan tiga digit menjadi ratusan, puluhan, dan satuan & 25 & CPMK0607\\
\barisgenap A2 & Tulis ekspresi logika untuk syarat yang diberikan dan tampilkan hasilnya & 25 & CPMK0607\\
B1 & Tunjukkan langkah evaluasi sebuah ekspresi campuran sesuai precedence & 20 & CPMK0608\\
\barisgenap B2 & Telusuri nilai variabel setelah penugasan majemuk dan increment & 20 & CPMK0608\\
B3 & Jelaskan beda \texttt{x++} dan \texttt{++x} dengan contoh & 10 & CPMK0608\\
\garisdasar
\end{tabular}
\vspace{10pt}

{\small Soal A dinilai dari keluaran yang tepat sama. Soal B dinilai dari ketepatan dan alasan. Pelanggaran aturan membuat nilai Exit Ticket ini 0.}
\note{Soal lengkap dibagikan terpisah dan tidak disimpan di repo publik.}
\end{frame}

\begin{frame}{Tugas 3: Kalkulator Lengkap dengan Validasi}
\framesubtitle{Dikumpulkan sebelum Pertemuan 4}
\begin{columns}[T]
\begin{column}{0.47\textwidth}
\begin{Blok}{Bagian A, program, 50 poin}
Buat program kalkulator yang membaca dua angka desimal, lalu menampilkan hasil berikut. Ketentuannya di slide berikut.
\end{Blok}
\end{column}
\begin{column}{0.47\textwidth}
\begin{BlokGelap}{Bagian B, penjelasan, 50 poin}
Komentar maksimal 150 kata: urutan evaluasi ekspresi syarat terakhir langkah demi langkah, alasan perlu cast sebelum \texttt{\%}, dan satu kesalahan precedence atau pembagian yang kalian temui.
\end{BlokGelap}
\end{column}
\end{columns}
\vspace{8pt}
{\small Data di tugas adalah data kalian sendiri. Rubrik lengkap ada di Modul Belajar 3.}
\note{Ingatkan bahwa penjelasan disalin dari AI mudah dikenali karena tidak menyebut pengalaman mereka sendiri.}
\end{frame}

\begin{frame}{Ketentuan Tugas 3}
\framesubtitle{Kalkulator Lengkap dengan Validasi}
\footnotesize
\begin{tabular}{@{}>{\raggedright\arraybackslash}p{87pt}>{\raggedright\arraybackslash}p{788pt}@{}}
\kepala{No} & \kepala{Ketentuan Bagian A}\\
1 & Masukan dua angka bertipe \texttt{double}\\
\barisgenap 2 & Hasil \texttt{+}, \texttt{-}, \texttt{*}, dan \texttt{/}\\
3 & Sisa bagi \texttt{\%} dari bagian bulat kedua angka (cast ke \texttt{int})\\
\barisgenap 4 & Sebelum membagi, periksa apakah pembagi nol; tampilkan pesan dengan operator ternary \texttt{?:}\\
5 & Nilai \texttt{bool}: apakah angka pertama lebih besar\\
\barisgenap 6 & Nilai \texttt{bool}: apakah hasil penjumlahan > 100 dan hasil kali > 50\\
Bonus & Tampilkan apakah bagian bulat masing-masing angka ganjil atau genap.\\
\garisdasar
\end{tabular}
\note{Bacakan ketentuan satu per satu dan tanyakan apakah ada yang belum jelas.}
\end{frame}

\begin{frame}{Rubrik Tugas 3}
\framesubtitle{Empat level, sama di semua instrumen}
\footnotesize
\begin{tabular}{@{}>{\raggedright\arraybackslash}p{166pt}>{\raggedright\arraybackslash}p{44pt}>{\raggedright\arraybackslash}p{166pt}>{\raggedright\arraybackslash}p{156pt}>{\raggedright\arraybackslash}p{146pt}>{\raggedright\arraybackslash}p{146pt}@{}}
\kepala{Kriteria} & \kepala{Poin} & \kepala{Sangat baik 85--100} & \kepala{Baik 70--84} & \kepala{Cukup 55--69} & \kepala{Kurang <55}\\
A1 Program berjalan & 20 & tanpa error dan peringatan & ada peringatan & perlu satu perbaikan & tidak terkompilasi\\
\barisgenap A2 Kelengkapan operasi & 15 & semua operasi dan validasi & satu operasi kurang & dua kurang & tiga atau lebih kurang\\
A3 Validasi dan logika & 15 & pembagi nol dan dua bool tepat & satu salah & dua salah & tidak ada validasi\\
\barisgenap B1 Langkah evaluasi & 30 & urutan tepat dan beralasan & satu langkah keliru & hanya hasil akhir & tidak ada atau disalin\\
B2 Cerita error dan pengujian & 20 & pesan, sebab, perbaikan, uji & pesan tidak dikutip & hanya perbaikan & tidak ada\\
\garisdasar
\end{tabular}
\vspace{8pt}

{\small Nilai kriteria = poin × skor level. Bagian A total 50, Bagian B total 50.}
\note{Rubrik ini sama strukturnya setiap minggu; yang berubah hanya isi kriteria A2, A3, dan B1.}
\end{frame}

\begin{frame}{Sebelum Pertemuan 4}
\framesubtitle{Persiapan}
\begin{columns}[T]
\begin{column}{0.31\textwidth}
\begin{Langkah}{1}{Selesaikan}
Hands-on yang belum lulus \texttt{cek.sh}, terutama latihan tantangan.
\end{Langkah}
\end{column}
\begin{column}{0.31\textwidth}
\begin{Langkah}{2}{Kumpulkan}
Tugas 3 sebelum Pertemuan 4 dimulai.
\end{Langkah}
\end{column}
\begin{column}{0.31\textwidth}
\begin{Langkah}{3}{Baca}
Modul Belajar 3 untuk mengulang, lalu bagian awal modul berikutnya.
\end{Langkah}
\end{column}
\end{columns}
\vspace{10pt}
Pertemuan 4 membahas percabangan dengan \texttt{if}, \texttt{else if}, dan \texttt{switch}, memakai ekspresi relasional dan logika dari minggu ini.\note{Tanyakan satu hal yang paling membingungkan hari ini untuk dibuka di review minggu depan.}
\end{frame}

\SlidePenutup{Terima kasih}{Sampai jumpa di Pertemuan 4: percabangan}
\note{Slide penutup.}
\end{document}
